\documentclass[a4paper]{article}
\input{Packages.tex}
\begin{document}


%%%%%%%%%%%%%%%%%%%%%%%%%%%%%%%%%%%%%%%%%%%%%%%%%%%%%%
%%                  0: title section                %%
%%%%%%%%%%%%%%%%%%%%%%%%%%%%%%%%%%%%%%%%%%%%%%%%%%%%%%

\title{Technical Note: \\
\emph{The Environmental Macroeconomics of the Circular Economy}}

\author{Rasmus Kehlet Skjødt Berg
% \thanks{University of Copenhagen, Oster Farimagsgade 5, 1353 Copenhagen K, Denmark. E-mail: rasmus.kehlet.berg@econ.ku.dk.}
\and
Peter Birch Sørensen
% \thanks{University of Copenhagen, Oster Farimagsgade 5, 1353 Copenhagen K, Denmark. E-mail: peter.birch.sorensen@econ.ku.dk.}
}
\maketitle

\begin{abstract}
\noindent
This note documents the three models used in the paper \emph{The Environmental Macroeconomics of the Circular Economy}: a planner's problem in continuous time, the corresponding decentralized market economy, and a quantitative implementation in discrete time with explicit functional forms. Its distinguishing feature is how the materials balance principle is imposed. Rather than positing a single aggregate identity relating extraction, capital accumulation and waste, conservation of mass is enforced separately for each economic activity --- production, extraction, exploration, waste collection and treatment, consumption, investment, and recycling --- and the aggregate relations used in all three models are derived from that process-level accounting rather than assumed directly.

Doing so makes clear that the materials balance is an \emph{accounting ledger} rather than a constraint on the allocation: it determines what the mass flows are, it does not restrict what the economy can do. The economy-wide waste-intensity level and the material intensity of the final good are accordingly endogenous objects, determined identically in the planner's problem and in the market. What the ledger leaves behind is the relation $W=R-\dot\Phi+\Omega\mathcal N$ --- waste generated equals material entering production, less the net accumulation of material embodied in the capital stock, plus mass displaced from nature that never entered production at all.

The last term is the model's one substantive fork, and the note is organized around it. The theory carries it throughout. The quantitative model sets it to zero, once and visibly, and that single restriction is what delivers the results the model is built to display: total waste becomes independent of output and consumption at the margin, so no Pigouvian tax on either is part of the optimal policy; the recycling margin prices recycled material at exactly the extraction cost, resource rent and pollution damage it saves, with the ledger price cancelling exactly; capital, because it stores material and releases it only as it depreciates, is a temporary materials sink whose social price is below unity, so that the instrument correcting the capital-formation margin is a \emph{subsidy} to gross investment rather than a tax; and the first-best is implemented with four instruments rather than the seven the general accounting requires. Each of these is a property of the restriction, not of materials-balance models in general, and the note says so at each point.
\end{abstract}


%%%%%%%%%%%%%%%%%%%%&
%%% Main sections %%%
%%%%%%%%%%%%%%%%%%%%%


\clearpage
\part{Models}\label{part:models}


This part contains a description of a number of related models. Section~\ref{sec:baseline} states which of them is the baseline and catalogues every point at which an alternative is entertained; the remaining sections present the three models in turn.

%%%%%%%%%%%%%%%%%%%%%%%%%%%%%%%%%%%%%%%%%%%%%%%%%%%%%%
%%  Part I, Section 0: The baseline model and its    %%
%%  switches                                         %%
%%%%%%%%%%%%%%%%%%%%%%%%%%%%%%%%%%%%%%%%%%%%%%%%%%%%%%

\section{The Baseline Model and Its Switches}
\label{sec:baseline}

This document states a single model three times --- as a planner's problem in continuous time (Section~\ref{sec:ac}), as the corresponding market economy (Section~\ref{sec:ad}), and as a discrete-time quantitative implementation with explicit functional forms (Section~\ref{sec:aq}) --- and along the way it considers a number of alternatives to particular modelling choices. This short section says which version is \emph{the} model, and catalogues every point at which an alternative is entertained. It exists so that no reader has to reconstruct the baseline by reading the remarks.

\subsection{The rule followed throughout}
\label{subsec:baseline:rule}

The main text always states the \textbf{most general version available}, and departures from it are set off explicitly rather than folded in. Concretely:

\begin{itemize}[leftmargin=1.6em,itemsep=0.25em]
\item Where a general case and a restricted case both exist and both are tractable, the general case is what the derivations use. The restricted case appears as a labelled special case, displayed alongside, never substituted in silently.
\item Where the quantitative model departs from the theory, it does so at \emph{one} identified point, Section~\ref{subsec:aq:B1}, and the departure is named and justified there.
\item Where two specifications are genuinely non-nested --- neither more general --- neither is privileged. Both are carried, and the choice between them is reported as a result rather than buried as an assumption.
\end{itemize}

Accordingly the switches below fall into three kinds, and it is worth keeping them distinct because they carry different burdens of proof.

\begin{description}[leftmargin=0em,itemsep=0.3em]
\item[\textnormal{\textbf{Restriction.}}] The main text carries the general case; the restriction is a special case obtained by setting parameters. Adopting it is a matter of convenience and its cost is computable. Every variant V\ref{var:aq:linear}--V\ref{var:aq:constant-phi} of Section~\ref{subsec:aq:variants} is of this kind, as is \textbf{(B1)}.
\item[\textnormal{\textbf{Fork.}}] Two non-nested specifications. Neither is more general, so the choice cannot be deferred to a robustness check --- it must be made, or both must be reported.
\item[\textnormal{\textbf{Maintained assumption.}}] A modelling choice whose alternative is described but \emph{not} derived anywhere in this document. These are the honest gaps: adopting the alternative would require work that has not been done, and the remarks say how much.
\end{description}

\subsection{The catalogue}
\label{subsec:baseline:catalogue}

\begin{table}[H]
\centering
\caption{Every modelling switch in this document, what the baseline sets it to, and where the alternative is discussed.}
\label{tab:baseline:switches}
\setlength{\tabcolsep}{4pt}
\renewcommand{\arraystretch}{1.05}
\begin{tabular}{@{}p{0.26\textwidth}p{0.09\textwidth}p{0.32\textwidth}p{0.25\textwidth}@{}}
\toprule
Switch & Kind & Baseline & Alternative, discussed in \\
\midrule
Nature-attributable waste, $\omega^{N,S},\omega^{D,S}$
& Restr.
& General in Sections~\ref{sec:ac}--\ref{sec:ad}; set to zero in Section~\ref{sec:aq} as \textbf{(B1)}
& Rem.~\ref{rem:ac:nature}, \ref{rem:aq:scope}; Sec.~\ref{subsec:aq:B1} \\
\addlinespace[0.2em]
Scaling of nature-attributable waste
& Fork
& Relative: $\mathcal M=\Omega\mathcal N$, scaled by the derived level $\Omega$
& Absolute intensities, Rem.~\ref{rem:ac:dilution} \\
\addlinespace[0.2em]
Closing rule for $\phi^I_t$
& Fork
& \emph{None declared.} Both (a) exogenous and (b) $\phi^I=\varsigma\Omega$ carried throughout
& Rem.~\ref{rem:ac:phiI}; Sec.~\ref{subsec:aq:phiI} \\
\addlinespace[0.2em]
Vintage structure of $\Phi$
& Restr.
& $\phi^I_t\neq\phi^K_t$; $\Phi$ a genuine state
& Constant intensity, V\ref{var:aq:constant-phi} \\
\addlinespace[0.2em]
Materials floor $\bar R$
& Restr.
& $\bar R>0$
& $\bar R=0$, V\ref{var:aq:no-floor} \\
\addlinespace[0.2em]
Pollution feedback on productivity
& Restr.
& $\gamma>0$
& $\gamma=0$, V\ref{var:aq:no-feedback} \\
\addlinespace[0.2em]
Endogenous reserves
& Restr.
& $S_t,X_t$ evolve; exploration active
& Fixed base, V\ref{var:aq:no-resource} \\
\addlinespace[0.2em]
Technical progress
& Restr.
& All $g^\cdot$ free
& All zero, V\ref{var:aq:no-progress} \\
\addlinespace[0.2em]
Recycling active
& Restr.
& Full model
& Linear economy, V\ref{var:aq:linear} \\
\addlinespace[0.2em]
Horizon
& Restr.
& Infinite in theory; finite $T$ for all numerical results
& V\ref{var:aq:finite} \\
\midrule
Timing of recycled input
& Maint.
& Contemporaneous: $R^R_t=a_t\varpi_tW_t$ within the period
& One-period lag, Rem.~\ref{rem:ac:recirculation}, \ref{rem:aq:period} \\
\addlinespace[0.2em]
Embodied-material pools
& Maint.
& One well-mixed stock $\Phi$; $K^Y$ and $K^R$ share $\phi^K$
& Two pools, Rem.~\ref{rem:ac:Phi-pools} \\
\addlinespace[0.2em]
Depreciated capital material
& Maint.
& Enters $W$ on the same footing as any other waste
& Costless re-embodiment, Rem.~\ref{rem:ac:depreciation} \\
\addlinespace[0.2em]
Curvature of $C^D$ in $D$
& Maint.
& Linear, giving a bang-bang exploration margin
& Strictly convex, Rem.~\ref{rem:aq:linear-costs} \\
\addlinespace[0.2em]
Instrument normalization
& Maint.
& $\tau^{W^R}=0$; output-side dilution loaded onto $\tau^Y$
& Rem.~\ref{rem:ad:normalization} \\
\bottomrule
\end{tabular}
\end{table}

\subsection{Reading the catalogue}
\label{subsec:baseline:reading}

Three comments, in decreasing order of importance.

\smalltitle{\textbf{(B1)} is the only restriction that separates the theory from the quantitative model.} Every other entry marked \emph{Restr.} is a variant used for building up, testing or decomposing the numerical implementation, and none of them is imposed in the headline results. \textbf{(B1)} is different: it is imposed throughout Section~\ref{sec:aq} and therefore in every number the model produces. It is singled out with its own label, its own subsection, and its own accounting of costs and benefits for exactly that reason. A reader who wants to know what the quantitative model assumes that the theory does not needs to read one thing, Section~\ref{subsec:aq:B1}.

\smalltitle{The two forks are genuinely open, and are treated as such.} Neither the closing rule for $\phi^I_t$ nor the scaling of nature-attributable waste is settled here. For $\phi^I_t$ this is deliberate and the quantitative section reports both rules; the two disagree about whether the relative intensities $\omega^i$ affect the allocation at all, which is too substantive a question to answer by picking a default. For the scaling, the document adopts the relative form for continuity with the process-level accounting of Section~\ref{subsec:ac:accounting}, but Remark~\ref{rem:ac:dilution} argues the absolute form is better identified and better behaved, and that the choice should be settled before the general accounting is ever taken to data. The relative form is a baseline by inheritance, not by argument.

\smalltitle{The maintained assumptions are where the unfinished work is.} Each of the five would require a derivation this document does not contain --- an extra state for the recycling lag, a second liability for two $\Phi$ pools, a re-derivation of $q$ for the depreciation bypass, the loss of a closed form for convex $C^D$. They are listed together so that the boundary of what has been established is visible in one place rather than distributed across five remarks. The first two are the ones most likely to matter quantitatively; Remark~\ref{rem:aq:period} explains why the recycling lag is not innocuous once the period length is fixed.

%%%%%%%%%%%%%%%%%%%%%%%%%%%%%%%%%%%%%%%%%%%%%%%%%%%%%%
%%  Alternative Part I, Section 1: Centralized theory %%
%%%%%%%%%%%%%%%%%%%%%%%%%%%%%%%%%%%%%%%%%%%%%%%%%%%%%%

\section{The Planner Model in Continuous Time}
\label{sec:ac}

This section states the continuous-time Ramsey model with exhaustible natural resources, waste generation, waste treatment and recycling, and characterizes the first-best allocation chosen by a utilitarian social planner. The section is self-contained: it states the primitives, states the planner's problem as an optimal control problem, lists the first-order and co-state conditions, and closes with a count of equations and unknowns and a set of technical remarks.

The distinguishing feature of the treatment below is how the materials balance principle is imposed. Conservation of mass is enforced separately for each economic activity --- production, extraction, exploration, waste collection and treatment, consumption, investment, and recycling --- and the aggregate waste-generation and materials-balance relations are then \emph{derived} by summing across activities, rather than posited directly as a single aggregate identity. This matters because an aggregate identity written down directly can look consistent while quietly mixing units (a final-good-denominated intensity applied to a physical tonnage, for instance) or silently assuming that some activities generate no waste at all; building the balance up from process-level mass conservation makes every such choice explicit and forces it to be stated as a modeling assumption rather than left implicit in a reduced-form equation.

A second feature follows from the first and shapes everything downstream. Conservation of mass is an \emph{accounting ledger}, not a technology: it does not restrict what the economy can do, it records what the mass flows are given what the economy did. The genuine restrictions are elsewhere and are already in the model --- reserves are finite, extraction is costly, recycling recovers only a fraction of what is fed into it, production requires materials. Accordingly, the material intensity of the final good, $\phi^Y$, and the economy-wide waste-intensity level, $\Omega$, are both \emph{endogenous} objects determined by the accounting, not calibrated parameters; what is calibrated is the vector of \emph{relative} intensities $\omega^i$ across activities. Section~\ref{subsec:ac:balances} makes this precise and Remark~\ref{rem:ac:ledger} explains why the alternative --- treating the balance as a constraint carrying its own shadow price --- is not available.

\smalltitle{Generality.} This section and Section~\ref{sec:ad} are stated under the \emph{general} accounting, in which extraction and exploration displace mass that never passed through the final good --- overburden and gangue --- at the nature-attributable intensities $\omega^{N,S}$ and $\omega^{D,S}$. This is the more general of the two settings catalogued in Section~\ref{sec:baseline}, and it is carried through every first-order condition, co-state condition and policy instrument below. The restriction $\omega^{N,S}=\omega^{D,S}=0$, which buys a closed-form ledger and makes the relative intensities allocation-irrelevant, is displayed at each point where it simplifies something --- always marked \textbf{(B1)} and set off from the general statement --- but it is not imposed here. Section~\ref{sec:aq} imposes it once and for all, for the reasons given there.

All variables are functions of time $t$; the time argument is suppressed unless it is needed. A dot denotes a time derivative.

\subsection{Notation}
\label{subsec:ac:notation}

Table~\ref{tab:ac:notation} collects the symbols used throughout this document. The same symbols are used in the decentralized model and, with a time subscript, in the quantitative model.

\begin{table}[H]
\centering
\caption{Notation used throughout this document.}
\label{tab:ac:notation}
\setlength{\tabcolsep}{4pt}
\renewcommand{\arraystretch}{0.75}
\begin{tabular}{@{}ll@{\hspace{1.5em}}ll@{}}
\toprule
\multicolumn{2}{@{}l}{\textbf{Stocks (state variables)}} & \multicolumn{2}{l@{}}{\textbf{Flows (controls and outcomes)}} \\
\midrule
$K$   & total man-made capital                & $C$   & consumption of final goods \\
$S$   & known reserves of natural resource    & $Y$   & output of final goods \\
$X$   & accumulated discoveries               & $I$   & net investment, $I=\dot K$ \\
$P$   & stock of pollution                    & $N$   & extraction of virgin materials \\
$\Phi$ & material embodied in capital stock   & $D$   & new reserves discovered \\
\multicolumn{2}{@{}l}{\textbf{Allocation of capital}}    & $R$   & total raw material input, $R=N+R^R$ \\
$K^Y$ & capital in final goods production     & $R^R$ & recycled materials \\
$K^R$ & capital in recycling                  & $W$   & total waste generated \\
      &                                       & $W^i$ & waste generated by activity $i$ \\
\multicolumn{2}{@{}l}{\textbf{Technology and taste functions}} & $\varpi$ & fraction of waste treated \\
$u(C)$        & utility of consumption        & $a$   & recycling rate \\
$v(P)$        & disutility of pollution       & $C^i$ & final-good cost of activity $i\in\{N,D,W\}$ \\
$F(K^Y,R,P)$  & final goods technology        & $\Xi$ & flow of waste reaching the environment \\
$a(\cdot)$    & recycling rate function       & \multicolumn{2}{l@{}}{\textbf{Shadow prices (units of the final good)}} \\
$C^N(N,S)$    & extraction cost               & $\lambda^K$ & marginal utility value of capital \\
$C^D(D,X)$    & discovery cost                & $\zeta$ & marginal utility value of the final good \\
$c^c$         & unit collection cost          & $p^S$ & marginal value of reserves \\
$c^T(\varpi)$ & unit treatment cost           & $p^X$ & marginal cost of accumulated discoveries \\
$\theta(P)$   & natural absorption rate       & $p^P$ & marginal cost of the pollution stock \\
                                             && $p^W$ & marginal cost of waste generation \\
\multicolumn{2}{@{}l}{\textbf{Derived accounting objects}} & $p^\Phi$ & liability of embodied capital material \\
$\Omega$   & economy-wide waste-intensity level & $B$ & social value of recycled material \\
$\phi^Y$   & material intensity of the final good & $q$ & social price of capital, $q=\lambda^K/\zeta$ \\
$\phi^K$   & average intensity of capital stock, $\Phi/K$ & $r$ & rate of return on capital \\
$\mathcal D$ & final-good-denominated waste base & $\iota$ & goods rate of interest \\
$\mathcal N$ & nature-attributable waste base    & $\pi$ & dilution credit, $\pi=p^W\Omega\kappa$ \\
$\mathcal M$ & nature-attributable waste flow, $\Omega\mathcal N$ & & \\
$\kappa$   & overburden ratio, $\mathcal N/\mathcal D$ & & \\
\midrule
\multicolumn{4}{@{}l@{}}{\textbf{Parameters:} $\rho$ (rate of time preference), $\delta$ (depreciation), $d^W$ (pollution potency of treated waste),}\\
\multicolumn{4}{@{}l@{}}{$\omega^Y,\omega^C$ (relative waste intensity of output net of depreciation, and of consumption),}\\
\multicolumn{4}{@{}l@{}}{$\omega^{N,Y},\omega^{N,S}$ (relative waste intensity of extraction, final-good- and nature-attributable),}\\
\multicolumn{4}{@{}l@{}}{$\omega^{D,Y},\omega^{D,S}$ (same, for exploration), $\omega^W$ (relative waste intensity of treatment),}\\
\multicolumn{4}{@{}l@{}}{$\phi^I(t)$ (material intensity of \emph{newly produced} capital),}\\
\multicolumn{4}{@{}l@{}}{$\varepsilon\equiv a'\tfrac{K^R/(\varpi W)}{a}$ (elasticity of the recycling rate).}\\
\bottomrule
\end{tabular}
\end{table}

Note what is \emph{absent} relative to a formulation in which the materials balance is imposed as a constraint: there is no multiplier $p^M$ on a materials balance, no composite $MSC^W$, and no waste-adjustment factor $\Lambda$ applied to marginal products. Section~\ref{subsec:ac:foc} shows why none of the three is needed.

\subsection{Tastes and technologies}
\label{subsec:ac:primitives}

\smalltitle{Preferences.} A representative family dynasty with an infinite horizon derives utility from consumption and disutility from the stock of pollution. Lifetime utility at time zero is
\begin{align}
U_0 &= \int_0^{\infty}\left[u(C)-v(P)\right]e^{-\rho t}\,dt,
&&
u'>0,\quad u''<0,\quad v'>0,\quad v''>0,\quad \rho>0.
\label{eq:ac:utility}
\end{align}

\smalltitle{Final goods.} Output is produced from capital and raw materials, and is reduced by pollution:
\begin{align}
Y &= F\left(K^Y,R,P\right),
&&
F_K>0,\ F_{KK}<0,\ F_R>0,\ F_{RR}<0,\ F_P<0.
\label{eq:ac:output}
\end{align}

\smalltitle{Materials and recycling.} Raw material input is the sum of virgin and recycled materials,
\begin{align}
R &= N + R^R,
\label{eq:ac:materials}
\end{align}
and recycled materials are produced from the treated share $\varpi W$ of waste using capital $K^R$:
\begin{align}
R^R &= a\!\left(x\right)\varpi W,
\qquad
x\equiv\frac{K^R}{\varpi W},
&&
a(0)=0,\quad a'>0,\quad a''<0,\quad
\varepsilon\equiv \frac{a'x}{a}<1,\quad
\lim_{x\to\infty}a=1.
\label{eq:ac:recycling}
\end{align}

\smalltitle{Allocation of capital.} The total capital stock is split between the two uses,
\begin{align}
K &= K^Y + K^R.
\label{eq:ac:capital-allocation}
\end{align}

\smalltitle{Extraction and exploration.} Virgin materials are extracted from a stock $S$ of known reserves, which is augmented by exploration:
\begin{align}
C^N &= C^N(N,S),
&&
C^N_N>0,\quad C^N_{NN}\ge 0,\quad C^N_S<0,
\label{eq:ac:extraction-cost}
\\
C^D &= C^D(D,X),
&&
C^D_D>0,\quad C^D_{DD}\ge 0,\quad C^D_X>0.
\label{eq:ac:discovery-cost}
\end{align}

\smalltitle{Waste collection and treatment.} All generated waste must be collected; the treated share $\varpi$ carries an increasing marginal treatment cost:
\begin{align}
C^W &= \left[c^c+c^T(\varpi)\right]W,
&&
c^T(0)=\left.\frac{dc^T}{d\varpi}\right|_{\varpi=0}=0,\quad
\frac{dc^T}{d\varpi}>0,\quad \frac{d^2c^T}{d\varpi^2}>0 \ \ \text{for }\varpi>0.
\label{eq:ac:waste-cost}
\end{align}

\smalltitle{Accumulation.} The aggregate resource constraint, the definition of investment, and the laws of motion of the three remaining stocks are
\begin{align}
Y &= C+I+\delta K+C^N+C^D+C^W,
&&
\dot K = I,
\label{eq:ac:resource-constraint}
\\
\dot S &= D-N,
&&
\dot X = D,
\label{eq:ac:reserves-motion}
\\
\dot P &= \Xi-\theta(P)P,
&&
0<d^W<1,\quad \theta'(P)<0,
\label{eq:ac:pollution-motion}
\end{align}
where the flow of waste that actually reaches the environment is
\begin{align}
\Xi &\equiv \left[1-\varpi+d^W\varpi\left(1-a\right)\right]W
\;=\; W\left(1-\varpi+d^W\varpi\right)-d^WR^R,
\label{eq:ac:Xi}
\end{align}
the second equality using $a\varpi W=R^R$ from~\eqref{eq:ac:recycling}. Writing $\Xi$ this way is not cosmetic: it is linear in $W$ and $R^R$ separately, so every derivative of the pollution flow below is immediate, and no derivative of $a(\cdot)$ with respect to $\varpi$ or $W$ ever needs to be taken. $W$ here is the total physical waste flow, whose determination --- the point of this section --- is worked out next.

\subsection{Materials accounting, process by process}
\label{subsec:ac:accounting}

Every process below is described by a mass balance: what enters equals what leaves, distributed across waste and embodied material. The final good $Y$ is abstract in the sense that it is a composite of many different physical activities, but it is produced from physical raw material and is itself spent on physical or quasi-physical activities (consumption, capital formation, extraction effort, and so on). Tracking the materials balance principle correctly requires following the mass embodied in $Y$ through each of those uses. $\Omega$ converts a flow of the final good into common waste units, and the relative intensities $\omega^i$ record how material-intensive the bundle drawn for purpose $i$ is compared with the bundle drawn for any other purpose.

Three modeling choices discipline the accounting throughout, stated here and motivated as each process is introduced below:
\begin{enumerate}[label=(\arabic*),leftmargin=2.4em]
\item Newly formed capital is embodied at the \emph{current} material intensity $\phi^I(t)$, not at the intensity of the capital it augments or replaces. The existing stock's average intensity $\phi^K(t)\equiv\Phi(t)/K(t)$ is therefore a genuine state variable, evolving as capital of one vintage is diluted or replaced by capital of another; Section~\ref{subsec:ac:embodied-stock} makes this precise. The case of a single constant intensity, $\phi^I(t)\equiv\phi^K(t)$ for all $t$, is nested as the special case in which the stock's composition has already converged.
\item Investment generates no waste of its own: all materials drawn for capital formation are embodied in the new capital, and waste associated with capital arises only through depreciation.
\item Waste generated by an activity that spends the final good is scaled by that activity's final-good \emph{cost}, not by the physical quantity it produces, whenever the two are measured in different units. This applies to extraction (scaled by $C^N$, not $N$), exploration (scaled by $C^D$, not $D$), and waste treatment (scaled by $C^W$, not $W$). Consumption needs no such correction, since $C$ is already measured in final-good units. Extraction and exploration additionally generate waste that is proportional to the physical flow itself ($N$, $D$) rather than to the resources spent obtaining it --- displaced material from the earth, brought up alongside the useful material, rather than a byproduct of processing the final good.
\end{enumerate}

\smalltitle{Production.} Production uses raw material $R=N+R^R$ and capital $K^Y$, which depreciates at rate $\delta$. Depreciation releases the capital's embodied material, $\phi^K\delta K^Y$, directly as waste; it is not a fresh draw on $R$:
\begin{align}
R+\phi^K\delta K^Y &= W^Y+M^Y, &
W^Y &= \Omega\omega^YY+\phi^K\delta K^Y, &
M^Y &= \phi^YY,
\label{eq:ac:production}
\end{align}
where $\phi^Y$ is the material intensity of the final good. The $\phi^K\delta K^Y$ terms cancel on both sides, leaving
\begin{align}
R &= \Omega\omega^YY+\phi^YY.
\label{eq:ac:production-reduced}
\end{align}
Producing output therefore requires no \emph{additional} raw material to cover capital depreciation: depreciation is a pure pass-through of previously embodied mass from the capital stock into the waste stream, and $\Omega\omega^Y$ measures the waste generated by production net of that pass-through.

\smalltitle{Extraction and exploration.} Extraction uses the final good, through the cost function $C^N(N,S)$, to pull $N$ tonnes from the reserve $S$. Two physically distinct sources of waste are at play: the embodied material of the final good spent on extraction effort itself (machinery, energy, and so on, all ultimately final-good expenditure), and displaced material from the earth --- overburden and gangue brought up alongside the ore --- proportional to the tonnage extracted rather than to the effort expended, and never part of $Y$ to begin with:
\begin{align}
W^N &= \Omega\omega^{N,Y}C^N(N,S)+\Omega\omega^{N,S}N.
\label{eq:ac:extraction}
\end{align}
The same logic applies to exploration, which spends $C^D(D,X)$ of the final good to add $D$ to the stock of discoveries:
\begin{align}
W^D &= \Omega\omega^{D,Y}C^D(D,X)+\Omega\omega^{D,S}D.
\label{eq:ac:exploration}
\end{align}
Only the final-good-attributable components, $\Omega\omega^{N,Y}C^N$ and $\Omega\omega^{D,Y}C^D$, are backed by embodied material drawn from $Y$; the nature-attributable components, $\Omega\omega^{N,S}N$ and $\Omega\omega^{D,S}D$, are additional mass that never passed through $Y$ and so never draws on $R$. This distinction turns out to be the \emph{only} channel through which the level $\Omega$ affects the allocation; see Remark~\ref{rem:ac:nature}.

\smalltitle{Waste collection and treatment, and consumption.} Both activities draw the final good and return its entire embodied content as waste, since neither stores materials durably: consumption goods are used up, and treating waste is itself an activity that consumes final good and leaves nothing durable behind. Consumption is already measured in final-good units, so $\Omega\omega^CC$ needs no further translation. Waste collection and treatment is not: the physical quantity it processes is $W$, a tonnage, while the final-good resources it uses are $C^W=\left[c^c+c^T(\varpi)\right]W$ from~\eqref{eq:ac:waste-cost}. As with extraction and exploration, it is the final-good expenditure that carries embodied material, so treatment waste is scaled by $C^W$ rather than by $W$ itself:
\begin{align}
W^W &= \Omega\omega^WC^W, & W^C &= \Omega\omega^CC.
\label{eq:ac:wc}
\end{align}

\smalltitle{Investment.} Gross resources devoted to capital formation are
\begin{align}
G &\equiv I+\delta K,
\label{eq:ac:gross-formation}
\end{align}
net accumulation plus replacement of depreciated capital, matching the $I$ and $\delta K$ terms in~\eqref{eq:ac:resource-constraint}. Under modeling choice~(2), none of this is wasted directly; under choice~(1), it is embodied at the \emph{current} intensity $\phi^I(t)$, since gross investment --- replacement no less than net accumulation --- is manufactured with today's technology, not the technology of whatever capital happens to be currently installed:
\begin{align}
W^I &= 0, & M^I &= \phi^IG.
\label{eq:ac:investment}
\end{align}

\smalltitle{Recycling.} Recycling uses capital $K^R$ and waste materials only, drawing no final good directly; its own depreciation releases waste in the same way as $K^Y$'s:
\begin{align}
W^{R,\text{waste}} &= \phi^K\delta K^R, & R^R &= a\!\left(x\right)\varpi W.
\label{eq:ac:recycling-mass}
\end{align}

\subsection{The embodied-material stock}
\label{subsec:ac:embodied-stock}

Because gross investment is embodied at the current intensity $\phi^I(t)$ while depreciation releases material at the stock's average intensity $\phi^K(t)$, the total mass of material currently embodied in the capital stock, $\Phi\equiv\phi^KK$, is itself a state variable:
\begin{align}
\dot\Phi &= \phi^I(t)\left(I+\delta K\right)-\delta\Phi,
&
\Phi(0) \text{ given},
\label{eq:ac:Phi-motion}
\end{align}
structurally identical to $K$'s own law of motion $\dot K=I$, with gross investment $I+\delta K$ replacing $I$ and the flow intensity $\phi^I(t)$ replacing the constant $1$. Equivalently, in ratio form,
\begin{align}
\dot\phi^K &= \left[\phi^I(t)-\phi^K\right]\frac{I+\delta K}{K},
\label{eq:ac:phiK-motion}
\end{align}
so $\phi^K(t)$ is an investment-weighted moving average of past values of $\phi^I$: it drifts toward the current vintage's intensity at a rate set by the gross investment rate, and is stationary exactly when $\phi^I(t)=\phi^K(t)$ --- the constant-intensity case of modeling choice~(1) is the special case in which the stock's composition has already converged.

$\Phi$ enters the rest of the accounting only through the depreciation pass-through terms already given in~\eqref{eq:ac:production} and~\eqref{eq:ac:recycling-mass}, $\phi^K\delta K^Y$ and $\phi^K\delta K^R$, which sum to $\phi^K\delta K=\delta\Phi$ under the maintained assumption that $K^Y$ and $K^R$ are drawn from the same, well-mixed stock and so share the same average intensity $\phi^K(t)$; see Remark~\ref{rem:ac:Phi-pools}. $\phi^I(t)$ itself is treated as a primitive of the model, on par with the $\omega^i$'s; Remark~\ref{rem:ac:phiI} discusses the two natural ways of specifying it.

\subsection{The materials ledger}
\label{subsec:ac:balances}

\smalltitle{Three relations.} Write the final-good-denominated waste base and the nature-attributable waste base as
\begin{align}
\mathcal D &\equiv \omega^YY+\omega^{N,Y}C^N(N,S)+\omega^{D,Y}C^D(D,X)+\omega^WC^W+\omega^CC,
&
\mathcal N &\equiv \omega^{N,S}N+\omega^{D,S}D.
\label{eq:ac:base}
\end{align}
$\mathcal D$ collects every flow that is measured in final-good units and therefore carries embodied material drawn through $R$; $\mathcal N$ collects the two physical flows that do not. Summing the process balances~\eqref{eq:ac:production}--\eqref{eq:ac:recycling-mass}, and using $\delta K^Y+\delta K^R=\delta K$ together with $\phi^K\delta K=\delta\Phi$, gives total waste generation:
\begin{align}
W &= \Omega\,\mathcal D+\Omega\,\mathcal N+\delta\Phi.
\label{eq:ac:waste-generation}
\end{align}
Conservation of the material embodied in $Y$ requires that every unit of it end up either as waste from one of the implicit downstream processes or as newly embodied capital, embodied at the current vintage's intensity:
\begin{align}
\phi^YY &= \Omega\left[\omega^{N,Y}C^N(N,S)+\omega^{D,Y}C^D(D,X)+\omega^WC^W+\omega^CC\right]+\phi^IG.
\label{eq:ac:consistency}
\end{align}
Substituting~\eqref{eq:ac:consistency} into the production balance~\eqref{eq:ac:production-reduced} eliminates $\phi^Y$ and gives the aggregate materials balance:
\begin{align}
N+R^R &= \Omega\,\mathcal D+\phi^IG.
\label{eq:ac:materials-balance}
\end{align}
The left-hand side nets out recycled material: a unit of $R^R$ substitutes one-for-one for a unit of virgin extraction $N$, since both enter production only through $R=N+R^R$. The right-hand side uses the \emph{gross} capital-formation flow $I+\delta K$ valued at the \emph{current} intensity $\phi^I$, because replacing depreciated capital draws fresh raw material just as any other use of $Y$ does, and the replacement units are built with today's technology. The nature-attributable waste terms do \emph{not} appear here, precisely because that waste was never backed by material drawn from $R$.

\smalltitle{What is endogenous.} Relations~\eqref{eq:ac:consistency} and~\eqref{eq:ac:materials-balance} are two equations in the two accounting objects $\phi^Y$ and $\Omega$, both of which are outcomes rather than parameters:
\begin{align}
\Omega &= \frac{N+R^R-\phi^IG}{\mathcal D},
&
\phi^Y &= \Omega\,\frac{\mathcal D-\omega^YY}{Y}+\phi^I\frac{G}{Y}.
\label{eq:ac:Omega}
\end{align}
$\phi^Y$ is the composition-weighted average material intensity of the $Y$ bundle: the final good is a composite of physically heterogeneous goods, the bundle drawn for consumption is not the same bundle as the one drawn for mining equipment, and $\phi^Y$ moves as the composition of the bundle moves. Treating it as a fixed parameter would tie $R/Y$ to a constant through~\eqref{eq:ac:production-reduced}, contradicting the substitutability of capital for materials in~\eqref{eq:ac:output} and defeating the purpose of having $R$ as an input at all. Neither~\eqref{eq:ac:consistency} nor~\eqref{eq:ac:materials-balance} therefore restricts the allocation: given any feasible path, the two determine $\Omega$ and $\phi^Y$ pointwise. Remark~\ref{rem:ac:ledger} discusses why this is the right reading and what it rules out.

\smalltitle{The ledger.} Substituting~\eqref{eq:ac:materials-balance} into~\eqref{eq:ac:waste-generation} cancels $\Omega$ and every $\omega^i$ from the first term, leaving
\begin{align}
\boxed{\;W \;=\; N+R^R-\phi^IG+\delta\Phi+\Omega\,\mathcal N \;=\; R-\dot\Phi+\mathcal M\;}
&&
\mathcal M &\equiv \Omega\,\mathcal N,
\label{eq:ac:ledger}
\end{align}
using~\eqref{eq:ac:Phi-motion} for the second equality. \emph{Waste generated equals the material entering production, less the net accumulation of material embodied in the capital stock, plus the mass displaced from nature that never entered production at all.} That is the entire physical content of the materials balance principle for this economy, and it is what the planner's problem below is constrained by. Note that $\phi^K$ appears nowhere in it.

The first two terms are the accounted flow: mass conservation pins waste down exactly, because every gram of it passed through $R$ on its way in. The third term is the unaccounted one, and it does not cancel, because that mass was never backed by $R$. It is the \emph{only} channel through which the level $\Omega$ --- and hence the calibration of the relative intensities $\omega^i$ --- touches the allocation.

\smalltitle{Solving the ledger.} Two things in~\eqref{eq:ac:ledger} depend on $W$ itself: recycled material, $R^R=a\varpi W$, and the waste base, through $C^W=\left[c^c+c^T(\varpi)\right]W$ inside $\mathcal D$. Split the base into the part free of $W$ and the part proportional to it,
\begin{align}
\mathcal D &= \mathcal D^0+\omega^Wc^WW,
&
\mathcal D^0 &\equiv \omega^YY+\omega^{N,Y}C^N+\omega^{D,Y}C^D+\omega^CC,
&
c^W &\equiv c^c+c^T(\varpi),
\label{eq:ac:base-split}
\end{align}
and note that $\mathcal N$ does not depend on $W$ at all. Substituting $\Omega=\left(R-\phi^IG\right)/\mathcal D$ from~\eqref{eq:ac:Omega} into~\eqref{eq:ac:ledger} and clearing the denominator turns the ledger into a \emph{quadratic} in $W$:
\begin{align}
\left(1-a\varpi\right)\omega^Wc^W\,W^2
+\Big[\left(1-a\varpi\right)\mathcal D^0-\left(N-\dot\Phi\right)\omega^Wc^W-a\varpi\,\mathcal N\Big]W
-\Big[\left(N-\dot\Phi\right)\mathcal D^0+\left(N-\phi^IG\right)\mathcal N\Big] &= 0.
\label{eq:ac:ledger-quadratic}
\end{align}
The leading coefficient is positive whenever $a\varpi<1$, and the constant term is negative whenever $N>\dot\Phi$ and $N>\phi^IG$ --- the latter being the feasibility restriction $\Omega\ge0$ of Remark~\ref{rem:ac:ledger} evaluated at $W=0$. A quadratic with a positive leading coefficient and a negative constant term has exactly one positive root, so \eqref{eq:ac:ledger-quadratic} determines $W$ uniquely and in closed form. Conservation of mass still fixes the waste flow pointwise; it simply no longer does so by a one-line division.

\smalltitle{(B1) The baseline restriction.} Setting the nature-attributable intensities to zero, $\mathcal N=0$, the constant term becomes $-\left(N-\dot\Phi\right)\mathcal D^0$ and the quadratic factors, with $W=\left(N-\dot\Phi\right)/\left(1-a\varpi\right)$ dividing out exactly and $\mathcal D^0$ dropping out entirely:
\begin{align}
W &= \frac{N-\dot\Phi}{1-a\varpi},
&
R &= \frac{N-a\varpi\dot\Phi}{1-a\varpi},
&
a\varpi&<1
&&\textbf{(B1)}.
\label{eq:ac:ledger-solved}
\end{align}
The factor $\left(1-a\varpi\right)^{-1}$ is the \emph{circularity multiplier}: the number of times a tonne of virgin material passes through production before it finally leaves the economy. When the capital stock is stationary in material terms ($\dot\Phi=0$), \eqref{eq:ac:ledger-solved} reads $N=W\left(1-a\varpi\right)$ --- virgin extraction equals waste not recycled --- and $N\to0$ as $a\varpi\to1$. A fully circular economy is one that extracts nothing.

That $\mathcal D^0$ vanishes from~\eqref{eq:ac:ledger-solved} is the precise sense in which the relative intensities are allocation-irrelevant under \textbf{(B1)}: the waste base is what $\Omega$ is measured \emph{against}, and with no unaccounted mass flow there is nothing for $\Omega$ to be applied \emph{to}.

\subsection{The planner's problem}
\label{subsec:ac:problem}

\begin{problem}[First best]\label{prob:ac:planner}
The planner chooses the controls $\left\{C,I,K^R,\varpi,N,D,W\right\}$ to maximize~\eqref{eq:ac:utility} subject to the state equations $\dot K=I$, \eqref{eq:ac:Phi-motion}, \eqref{eq:ac:reserves-motion} and~\eqref{eq:ac:pollution-motion}, the static resource constraint obtained by substituting~\eqref{eq:ac:output}--\eqref{eq:ac:waste-cost} into~\eqref{eq:ac:resource-constraint} with $\dot K$ replaced by $I$,
\begin{align}
I &= F\!\left(K-K^R,\ N+a\!\left(x\right)\varpi W,\ P\right)-\delta K-C-C^N(N,S)-C^D(D,X)-\left[c^c+c^T(\varpi)\right]W,
\label{eq:ac:capital-constraint}
\end{align}
and the materials ledger~\eqref{eq:ac:ledger}, given the predetermined initial stocks $K(0),\Phi(0),S(0),X(0),P(0)$ and the non-negativity restrictions $K^R\ge 0$, $\varpi\in[0,1]$, $W\ge 0$.
\end{problem}

$W$ is listed as a control for convenience only: the ledger pins it down through~\eqref{eq:ac:ledger-quadratic}, and the first-order condition for $W$ therefore determines the associated multiplier rather than a behavioral margin. The level $\Omega$ is \emph{not} a free control either; it is the ratio~\eqref{eq:ac:Omega}, and it enters the problem only through the product $\Omega\mathcal N$ in the ledger. Under \textbf{(B1)} it drops out of Problem~\ref{prob:ac:planner} altogether and is recovered from~\eqref{eq:ac:Omega} after the allocation is known.

\smalltitle{Hamiltonian and multipliers.} Let $\lambda^K,\lambda^\Phi,\lambda^S,\lambda^X,\lambda^P$ denote the current-value co-state variables of the five stocks, $\zeta$ the multiplier on~\eqref{eq:ac:capital-constraint}, and $\eta$ the multiplier on the ledger~\eqref{eq:ac:ledger}, all in units of utility. The current-value Hamiltonian is
\begin{align}
\mathcal H =\;& u(C)-v(P)+\lambda^KI+\lambda^\Phi\left[\phi^IG-\delta\Phi\right]+\lambda^S\left(D-N\right)+\lambda^XD+\lambda^P\left[\Xi-\theta(P)P\right]
\notag \\
&+\zeta\Big[F\!\left(K-K^R,N+R^R,P\right)-\delta K-C-C^N(N,S)-C^D(D,X)-\left(c^c+c^T(\varpi)\right)W-I\Big]
\notag \\
&+\eta\Big[N+R^R-\phi^IG+\delta\Phi-W+\Omega\,\mathcal N\Big],
\qquad
\Omega=\frac{N+R^R-\phi^IG}{\mathcal D}.
\label{eq:ac:hamiltonian}
\end{align}
All prices below are expressed in units of the final good by normalizing with $\zeta$, with the sign chosen so that the price of every \emph{bad} is positive:
\begin{align}
p^S &\equiv \frac{\lambda^S}{\zeta},
&
p^X &\equiv -\frac{\lambda^X}{\zeta},
&
p^P &\equiv -\frac{\lambda^P}{\zeta},
&
p^\Phi &\equiv -\frac{\lambda^\Phi}{\zeta},
&
p^W &\equiv -\frac{\eta}{\zeta}.
\label{eq:ac:normalizations}
\end{align}
$p^S,p^X,p^P$ and $p^\Phi$ are shown to be positive in Section~\ref{subsec:ac:costates}. The sign of $p^W$ is \emph{not} assumed: it is an outcome, discussed in Remark~\ref{rem:ac:pW-sign}.

\smalltitle{Derivatives.} With $z\equiv\varpi W$ so that $R^R=a(K^R/z)z$,
\begin{align}
\frac{\partial R^R}{\partial K^R}&=a'(x),
&
\frac{\partial R^R}{\partial\varpi}&=a\left(1-\varepsilon\right)W,
&
\frac{\partial R^R}{\partial W}&=a\left(1-\varepsilon\right)\varpi,
\label{eq:ac:RR-derivatives}
\end{align}
and, from the linear form of $\Xi$ in~\eqref{eq:ac:Xi},
\begin{align}
\frac{\partial\Xi}{\partial K^R}&=-d^Wa'(x),
&
\frac{\partial\Xi}{\partial\varpi}&=-\left(1-d^W\right)W-d^Wa\left(1-\varepsilon\right)W,
&
\frac{\partial\Xi}{\partial W}&=\left(1-\varpi+d^W\varpi\right)-d^Wa\left(1-\varepsilon\right)\varpi.
\label{eq:ac:Xi-derivatives}
\end{align}

\smalltitle{How the nature-attributable term enters.} The ledger bracket in~\eqref{eq:ac:hamiltonian} differs from its \textbf{(B1)} counterpart only by $\Omega\mathcal N$, and because $\Omega$ is the ratio~\eqref{eq:ac:Omega} its derivative decomposes into three pieces. For any control $\upsilon$,
\begin{align}
\frac{\partial\left(\Omega\mathcal N\right)}{\partial\upsilon}
&= \underbrace{\kappa\,\frac{\partial\left(R-\phi^IG\right)}{\partial\upsilon}}_{\text{scale}}
+\underbrace{\Omega\,\frac{\partial\mathcal N}{\partial\upsilon}}_{\text{displacement}}
-\underbrace{\Omega\kappa\,\frac{\partial\mathcal D}{\partial\upsilon}}_{\text{dilution}},
&
\kappa &\equiv \frac{\mathcal N}{\mathcal D}.
\label{eq:ac:OmegaN-derivative}
\end{align}
Multiplying by $\eta=-\zeta p^W$ and dividing by $\zeta$, every first-order condition of Section~\ref{subsec:ac:foc} therefore carries, relative to its \textbf{(B1)} form, the three extra terms
\begin{align}
-\,p^W\kappa\,\frac{\partial\left(R-\phi^IG\right)}{\partial\upsilon}
\;-\;p^W\Omega\,\frac{\partial\mathcal N}{\partial\upsilon}
\;+\;\pi\,\frac{\partial\mathcal D}{\partial\upsilon},
&&
\pi &\equiv p^W\Omega\kappa.
\label{eq:ac:extra-terms}
\end{align}
The three have distinct readings. The \emph{scale} term says that a control which draws more mass into production raises nature-attributable waste in proportion, so the ledger price is effectively $p^W(1+\kappa)$ rather than $p^W$ wherever mass enters. The \emph{displacement} term is the direct one: extracting a tonne displaces $\Omega\omega^{N,S}$ tonnes of overburden. The \emph{dilution} term is the awkward one --- spending more final good on any activity in $\mathcal D$ enlarges the base against which $\Omega$ is measured, lowering $\Omega$ and hence lowering $\mathcal M$ --- and $\pi$ is the resulting credit per unit of final good added to the base. Remark~\ref{rem:ac:dilution} argues that this channel is an artifact of the specification rather than economics, and says what to do about it.

The derivatives needed are collected once here. Writing $c^W\equiv c^c+c^T(\varpi)$,
\begin{align}
\frac{\partial\mathcal D}{\partial C}&=\omega^C,
&
\frac{\partial\mathcal D}{\partial N}&=\omega^YF_R+\omega^{N,Y}C^N_N,
&
\frac{\partial\mathcal D}{\partial D}&=\omega^{D,Y}C^D_D,
&
\frac{\partial\mathcal D}{\partial K^R}&=\omega^Y\left[F_Ra'-F_K\right],
\label{eq:ac:D-derivatives-a}
\\
\frac{\partial\mathcal D}{\partial\varpi}&=\left[\omega^YF_Ra\left(1-\varepsilon\right)+\omega^W\frac{dc^T}{d\varpi}\right]W,
&
\frac{\partial\mathcal D}{\partial W}&=\omega^YF_Ra\left(1-\varepsilon\right)\varpi+\omega^Wc^W,
&
\frac{\partial\mathcal D}{\partial I}&=0,
\label{eq:ac:D-derivatives-b}
\end{align}
with $\partial\mathcal N/\partial N=\omega^{N,S}$, $\partial\mathcal N/\partial D=\omega^{D,S}$ and $\partial\mathcal N/\partial\upsilon=0$ for every other control, and with $\partial\left(R-\phi^IG\right)/\partial\upsilon$ equal to $1$, $-\phi^I$, $a'$, $a(1-\varepsilon)W$ and $a(1-\varepsilon)\varpi$ for $\upsilon=N,I,K^R,\varpi,W$ respectively and zero for $C$ and $D$.

\subsection{First-order conditions}
\label{subsec:ac:foc}

Each condition is stated in general form first, with its \textbf{(B1)} collapse displayed alongside. Throughout, $\kappa=\mathcal N/\mathcal D$ and $\pi=p^W\Omega\kappa$, and both vanish under \textbf{(B1)}.

\smalltitle{Optimal saving.} $\partial\mathcal H/\partial C=u'(C)-\zeta+\eta\left(-\Omega\kappa\,\omega^C\right)=0$:
\begin{align}
u'(C) &= \zeta\left[1-\pi\,\omega^C\right]
&&\xrightarrow{\ \textbf{(B1)}\ }
&
u'(C) &= \zeta.
\label{eq:ac:foc-saving}
\end{align}
Under \textbf{(B1)} there is no wedge between the marginal utility of consumption and the marginal utility value of the final good: consumption generates no waste \emph{at the margin}, because by~\eqref{eq:ac:ledger} $W$ then does not depend on $C$, and the mass embodied in consumption goods was already counted when it entered the economy as $R$. In general a wedge survives, but it is a \emph{dilution} wedge and not a waste-generation one: more consumption enlarges $\mathcal D$, which lowers $\Omega$, which lowers the overburden displaced per tonne extracted. With $p^W>0$ this makes $u'(C)<\zeta$ --- consumption is favoured at the margin. See Remark~\ref{rem:ac:dilution}.

\smalltitle{Optimal investment.} $\partial\mathcal H/\partial I=\lambda^K+\lambda^\Phi\phi^I-\zeta-\eta\phi^I\left(1+\kappa\right)=0$, i.e.
\begin{align}
q &\equiv \frac{\lambda^K}{\zeta}=1-\phi^I\left[p^W\left(1+\kappa\right)-p^\Phi\right]
&&\xrightarrow{\ \textbf{(B1)}\ }
&
q &= 1-\phi^I\left(p^W-p^\Phi\right)\;<\;1.
\label{eq:ac:foc-investment}
\end{align}
$\zeta$ is the marginal utility value of a unit of the \emph{final good} and $\lambda^K$ that of a unit of \emph{capital}; they differ by the factor $q$, the social price of capital in units of the final good. That price is strictly \emph{below} one: a unit of new capital removes $\phi^I$ tonnes of material from today's waste stream and stores them, and although the stored mass is released later as the capital depreciates --- the liability $p^\Phi$ --- the release is spread over time and discounted, so $p^\Phi<p^W$ always (Section~\ref{subsec:ac:costates}). \textbf{Capital is a temporary materials sink, and the sink is socially valuable.} The general form says the sink is worth \emph{more} than under \textbf{(B1)}, by the factor $\left(1+\kappa\right)$ on $p^W$: material stored in capital is material not drawn through production, and each tonne drawn through production drags $\kappa$ tonnes of overburden with it.

\smalltitle{Optimal extraction.} $\partial\mathcal H/\partial N=0$ gives, using~\eqref{eq:ac:extra-terms},
\begin{align}
F_R-C^N_N-p^W\left(1+\kappa\right)-p^W\Omega\,\omega^{N,S}+\pi\left[\omega^YF_R+\omega^{N,Y}C^N_N\right] &= p^S
\label{eq:ac:foc-extraction}
\\
\xrightarrow{\ \textbf{(B1)}\ }\qquad
F_R-C^N_N-p^W &= p^S.
\notag
\end{align}
A marginal tonne of virgin material is productive ($F_R$), costs $C^N_N$ to extract, depletes the reserve ($p^S$), and --- by the ledger --- adds one tonne to the waste stream, plus $\kappa$ tonnes through the scale channel and $\Omega\omega^{N,S}$ tonnes of displaced overburden, less whatever its own extraction cost dilutes away.

\smalltitle{Optimal exploration.} $\partial\mathcal H/\partial D=0$:
\begin{align}
C^D_D\left[1-\pi\,\omega^{D,Y}\right]+p^X+p^W\Omega\,\omega^{D,S} &= p^S
&&\xrightarrow{\ \textbf{(B1)}\ }
&
C^D_D+p^X &= p^S.
\label{eq:ac:foc-exploration}
\end{align}
Exploration draws no material through production, so it carries no scale term; it enters the general ledger only by displacing mass directly ($\omega^{D,S}$) and by diluting the base ($\omega^{D,Y}$).

\smalltitle{Optimal investment in recycling.} Using~\eqref{eq:ac:RR-derivatives}, \eqref{eq:ac:Xi-derivatives} and~\eqref{eq:ac:extra-terms}, $\partial\mathcal H/\partial K^R=0$ collects into a form identical to the \textbf{(B1)} one once two objects are defined --- the social value of a marginal unit of recycled material and the social opportunity cost of a unit of capital moved into recycling:
\begin{align}
B &\equiv F_R\left[1+\pi\,\omega^Y\right]-p^W\left(1+\kappa\right)+d^Wp^P,
&
\widehat F_K &\equiv F_K\left[1+\pi\,\omega^Y\right],
\label{eq:ac:B}
\end{align}
both of which collapse to $B=F_R-p^W+d^Wp^P$ and $\widehat F_K=F_K$ under \textbf{(B1)}. Imposing $K^R\ge0$, the condition takes complementary-slackness form:
\begin{subequations}\label{eq:ac:foc-recycling}
\begin{align}
K^R&=0
&&\text{if}\quad a'(0)\,B\le \widehat F_K,
\label{eq:ac:foc-recycling-a}
\\
a'(x)\,B&=\widehat F_K
&&\text{if}\quad a'(0)\,B> \widehat F_K.
\label{eq:ac:foc-recycling-b}
\end{align}
\end{subequations}
The factor $\left[1+\pi\omega^Y\right]$ appears on both sides because moving capital into recycling changes output, and output is itself in the waste base $\mathcal D$: the forgone $F_K$ shrinks the base and so raises $\Omega$, while the recycled material recovered enlarges it again through $F_Ra'$. It is common to $B$ and $\widehat F_K$ and therefore cancels from the ratio, which is why the closed forms of Section~\ref{subsec:aq:recycling-properties} survive intact.

\smalltitle{Optimal waste treatment.} Dividing $\partial\mathcal H/\partial\varpi=0$ by $\zeta W$:
\begin{align}
a\left(1-\varepsilon\right)B+p^P\left(1-d^W\right) &= \frac{dc^T}{d\varpi}\left[1-\pi\,\omega^W\right]
&&\xrightarrow{\ \textbf{(B1)}\ }
&
a\left(1-\varepsilon\right)B+p^P\left(1-d^W\right) &= \frac{dc^T}{d\varpi}.
\label{eq:ac:foc-treatment}
\end{align}
Treating a marginal tonne yields $a(1-\varepsilon)$ tonnes of recycled material, worth $B$ each, and reduces the flow reaching the environment by $\left(1-d^W\right)$; in general the marginal treatment cost is credited by $\pi\omega^W$ for the base it adds. Condition~\eqref{eq:ac:foc-treatment} is stated as an interior condition; see Remark~\ref{rem:ac:varpi} for the bounds on $\varpi$.

\smalltitle{The multiplier on the ledger.} Dividing $\partial\mathcal H/\partial W=0$ by $\zeta$:
\begin{align}
\varpi\,a\left(1-\varepsilon\right)B+p^W &= \left[c^c+c^T(\varpi)\right]\left[1-\pi\,\omega^W\right]+p^P\left(1-\varpi+d^W\varpi\right).
\label{eq:ac:foc-waste}
\end{align}
This is not a behavioral margin --- $W$ is pinned by~\eqref{eq:ac:ledger-quadratic} --- but the equation that determines $p^W$, which is what keeps the system square.

\subsection{Two reductions}
\label{subsec:ac:reductions}

The first-order conditions simplify considerably once the extraction condition is used to eliminate $F_R$, and the resulting forms are both more interpretable and better conditioned numerically.

\smalltitle{The value of recycled material.} Substituting the extraction condition~\eqref{eq:ac:foc-extraction} into~\eqref{eq:ac:B}, the terms in $F_R$ --- including the whole dilution factor $\left[1+\pi\omega^Y\right]$ --- cancel identically, leaving
\begin{align}
\boxed{\;B \;=\; \bar B+p^W\Upsilon\;}
&&
\bar B &\equiv p^S+C^N_N+d^Wp^P,
&
\Upsilon &\equiv \Omega\left[\omega^{N,S}-\kappa\,\omega^{N,Y}C^N_N\right],
\label{eq:ac:B-reduced}
\end{align}
and under \textbf{(B1)}, where $\Upsilon=0$,
\begin{align}
B &= p^S+C^N_N+d^Wp^P
&&\textbf{(B1)}.
\label{eq:ac:B-reduced-B1}
\end{align}
\emph{The marginal social value of recycled material is the extraction cost it saves, plus the reserve rent it saves, plus the pollution it avoids --- and, in general, whatever the tonne of virgin extraction it displaces would have dragged out of the ground alongside it.} That the ledger price $p^W$ enters only through $\Upsilon$ is the sharpest available check on the accounting: a unit of recycled material displaces a unit of virgin extraction one for one, so the two change total \emph{accounted} waste identically at the margin and $p^W$ can survive only on the \emph{unaccounted} channel. Note that $\Upsilon$ involves only the two extraction intensities: if extraction generates no waste of its own, $\omega^{N,S}=\omega^{N,Y}=0$, then $p^W$ cancels from $B$ exactly even when exploration displaces mass ($\omega^{D,S}>0$) and $\kappa\neq0$.

In particular, under \textbf{(B1)} the switch point in~\eqref{eq:ac:foc-recycling-a} is a comparison of $a'(0)\left[p^S+C^N_N+d^Wp^P\right]$ with $F_K$: the date at which the economy turns circular depends on reserve scarcity, extraction cost and pollution damage, and on none of the waste-intensity parameters. In general the waste-intensity parameters re-enter it, through $\Upsilon$ and through the common factor $\left[1+\pi\omega^Y\right]$ --- which cancels from the comparison, since it multiplies both sides.

\smalltitle{The ledger price.} Rearranging~\eqref{eq:ac:foc-waste} before substituting shows that $p^W$ carries the marginal circularity multiplier $\left[1-\varpi a(1-\varepsilon)\right]^{-1}$ --- the counterpart of $\left(1-a\varpi\right)^{-1}$ in~\eqref{eq:ac:ledger-solved}, reduced by the dilution of a fixed $K^R$ over a larger flow of waste. Substituting~\eqref{eq:ac:B-reduced} makes the multiplier cancel against the extraction margin. Because $B$ and $\pi$ are both linear in $p^W$, what remains is linear in $p^W$ and solves explicitly:
\begin{align}
\boxed{\;
p^W = \frac{c^c+c^T(\varpi)+p^P\left(1-\varpi+d^W\varpi\right)-\varpi\,a\left(1-\varepsilon\right)\bar B}
{1+\Omega\kappa\,\omega^W\left[c^c+c^T(\varpi)\right]+\varpi\,a\left(1-\varepsilon\right)\Upsilon}\;}
\label{eq:ac:pW-explicit}
\end{align}
The numerator is exactly the \textbf{(B1)} ledger price, and the denominator is unity under \textbf{(B1)}:
\begin{align}
p^W &= c^c+c^T(\varpi)+p^P\left(1-\varpi+d^W\varpi\right)-\varpi\,a\left(1-\varepsilon\right)\left[p^S+C^N_N+d^Wp^P\right]
&&\textbf{(B1)}.
\label{eq:ac:pW-explicit-B1}
\end{align}
Collection and treatment cost, plus pollution damage, less the recycling value of the tonne. In the linear stage ($a=0,\varpi=0$) this is simply $p^W=c^c+p^P$ in both cases. The general accounting therefore rescales the ledger price but does not change its composition --- a fact worth exploiting numerically, since it means the \textbf{(B1)} solver can be reused with a single correction factor rather than rewritten.

\subsection{Co-state conditions}
\label{subsec:ac:costates}

\smalltitle{Capital, and the price of capital.} $\Phi$'s law of motion and the ledger both contain $\phi^I\delta K$, so $K$ enters through them as well as through production:
\begin{align}
\frac{\partial\mathcal H}{\partial K} &= \lambda^\Phi\phi^I\delta+\zeta\left(F_K-\delta\right)-\eta\left[\phi^I\delta\left(1+\kappa\right)+\Omega\kappa\,\omega^YF_K\right]
= \zeta\left[\widehat F_K-\delta q\right],
\label{eq:ac:H-K}
\end{align}
using~\eqref{eq:ac:foc-investment} and the definition of $\widehat F_K=F_K\left[1+\pi\omega^Y\right]$ from~\eqref{eq:ac:B}. \textbf{Depreciation is valued at $q$, not at $1$}, because replacing worn capital draws fresh material through the ledger and re-embodies it in the stock. Hence $\rho\lambda^K-\dot\lambda^K=\zeta\left[\widehat F_K-\delta q\right]=\lambda^Kr$ with
\begin{align}
r &\equiv \frac{\widehat F_K}{q}-\delta
\ \ \xrightarrow{\ \textbf{(B1)}\ }\ \ \frac{F_K}{q}-\delta,
&
\lambda^K(t) &= \lambda^K(0)\exp\left(-\int_0^t\left[r(u)-\rho\right]du\right).
\label{eq:ac:r}
\end{align}
A unit of capital earns the marginal product $\widehat F_K$ --- its marginal product in production, plus the dilution credit that output carries in the general accounting --- on the $q$ units of final good it costs, net of depreciation: the effective return is $r$.

\smalltitle{The goods rate of interest.} From $\lambda^K=\zeta q$ and~\eqref{eq:ac:r},
\begin{align}
\frac{\dot\zeta}{\zeta} &= \rho-\iota,
&
\iota &\equiv r+\frac{\dot q}{q}=\frac{\widehat F_K}{q}-\delta+\frac{\dot q}{q}.
\label{eq:ac:iota}
\end{align}
$\iota$ is the rate at which the planner trades \emph{goods} across time: the return on capital plus the capital gain or loss on the materials capital embodies. It, not $r$, is the discount rate for every price normalized by $\zeta$. The two coincide only when $q$ is constant, and both reduce to $F_K-\delta$ when newly formed capital embodies no materials ($\phi^I=0$).

\smalltitle{The three resource and environment stocks.} With $\partial\mathcal H/\partial S=-\zeta C^N_S$, $\partial\mathcal H/\partial X=-\zeta C^D_X$ and $\partial\mathcal H/\partial P=-v'(P)+\zeta F_P+\zeta p^P\widehat\theta$, the normalized co-state equations are $\dot p^S=\iota p^S+C^N_S$, $\dot p^X=\iota p^X-C^D_X$ and $\dot p^P=\left(\iota+\widehat\theta\right)p^P-MSC^P$, which integrate to
\begin{align}
p^S(t) &= \int_t^{\infty}\left(-C^N_S(z)\right)\exp\left(-\int_t^{z}\iota(u)\,du\right)dz\;>\;0,
\label{eq:ac:pS}
\\
p^X(t) &= \int_t^{\infty}C^D_X(z)\exp\left(-\int_t^{z}\iota(u)\,du\right)dz\;>\;0,
\label{eq:ac:pX}
\\
p^P(t) &= \int_t^{\infty}MSC^P(z)\exp\left(-\int_t^{z}\left[\iota(u)+\widehat\theta(u)\right]du\right)dz\;>\;0,
\label{eq:ac:pP}
\end{align}
where
\begin{align}
MSC^P &\equiv \frac{v'(P)}{\zeta}-F_P\;>\;0,
&
\widehat\theta &\equiv \frac{d\left[\theta(P)P\right]}{dP}=\theta\left(1-\varepsilon^P\right),
&
\varepsilon^P &\equiv -\theta'(P)\frac{P}{\theta}.
\label{eq:ac:MSCP}
\end{align}

\smalltitle{Embodied capital material.} $\Phi$ enters only through the depreciation pass-through $\delta\Phi$ in the ledger~\eqref{eq:ac:ledger} and through its own decay in~\eqref{eq:ac:Phi-motion}:
\begin{align}
\frac{\partial\mathcal H}{\partial\Phi} &= -\lambda^\Phi\delta+\eta\delta=\zeta\delta\left(p^\Phi-p^W\right),
\label{eq:ac:H-Phi}
\end{align}
so $\dot p^\Phi=\left(\iota+\delta\right)p^\Phi-\delta p^W$, which integrates to
\begin{align}
p^\Phi(t) &= \int_t^{\infty}\delta\,p^W(z)\exp\left(-\int_t^{z}\left[\iota(u)+\delta(u)\right]du\right)dz.
\label{eq:ac:pPhi}
\end{align}
$p^\Phi$ is the \emph{liability} attached to a unit of material already embodied in the capital stock: the present value of the waste it will generate as it depreciates, discounted at the goods rate of interest plus its own decay rate. Structurally this is~\eqref{eq:ac:pP} with $\delta$ in place of $\widehat\theta$ and $\delta p^W$ in place of $MSC^P$.

Comparing~\eqref{eq:ac:pPhi} with $p^W$ itself establishes the claim used in~\eqref{eq:ac:foc-investment}. Were $p^W$ constant, \eqref{eq:ac:pPhi} would give $p^\Phi=\delta p^W/\left(\iota+\delta\right)<p^W$; the same comparison holds for any positive path of $p^W$, since $\int_t^\infty\delta e^{-\int_t^z(\iota+\delta)}dz<1$ whenever $\iota>0$. Hence
\begin{align}
q &= 1-\phi^Ip^W\left[\kappa+\frac{\iota}{\iota+\delta}\right]
\ \ \xrightarrow{\ \textbf{(B1)}\ }\ \
1-\phi^Ip^W\,\frac{\iota}{\iota+\delta}\;<\;1
&&\text{(constant-}p^W\text{ case)},
\label{eq:ac:q-explicit}
\end{align}
so the materials value of capital is the value of \emph{deferring} a tonne of waste by embodying it in a stock that decays at rate $\delta$ --- worth the fraction $\iota/(\iota+\delta)$ of the tonne --- plus, in general, the value of the $\kappa$ tonnes of overburden that the deferred tonne does not drag out of the ground in the meantime. Since $\kappa\ge0$, $q$ is unambiguously further below one under the general accounting than under \textbf{(B1)}. If $p^W<0$ (Remark~\ref{rem:ac:pW-sign}) both $p^\Phi$ and $q-1$ change sign together.

\smalltitle{Transversality.} All five stocks require a terminal condition:
\begin{align}
\lim_{t\to\infty}\lambda^KK &= \lim_{t\to\infty}\lambda^\Phi\Phi = \lim_{t\to\infty}\lambda^SS = \lim_{t\to\infty}\lambda^XX = \lim_{t\to\infty}\lambda^PP = 0.
\label{eq:ac:transversality}
\end{align}

\smalltitle{Keynes--Ramsey rule.} Differentiating~\eqref{eq:ac:foc-saving} and using~\eqref{eq:ac:iota}:
\begin{align}
\frac{\dot C}{C} &= \frac{1}{\sigma}\left[\iota-\rho\right],
&
\sigma&\equiv-\frac{u''C}{u'}>0.
\label{eq:ac:keynes-ramsey}
\end{align}
The textbook rule, with the goods rate of interest $\iota$ in place of the net marginal product of capital. All of the model's materials machinery enters consumption growth through $\iota$ and through nothing else.

\subsection{The system: unknowns, equations, and stages}
\label{subsec:ac:system}

\smalltitle{Counting.} Given the states $\left(K,\Phi,S,X,P\right)$, the conditions above form a system of ten equations --- \eqref{eq:ac:foc-investment}, \eqref{eq:ac:foc-extraction}, \eqref{eq:ac:foc-exploration}, \eqref{eq:ac:foc-recycling}, \eqref{eq:ac:foc-treatment}, \eqref{eq:ac:pW-explicit}, \eqref{eq:ac:ledger-quadratic}, \eqref{eq:ac:Omega}, \eqref{eq:ac:capital-constraint} and the Keynes--Ramsey rule~\eqref{eq:ac:keynes-ramsey} --- in the ten unknowns
\begin{align*}
C,\quad I,\quad x,\quad \varpi,\quad N,\quad D,\quad W,\quad p^W,\quad q,\quad \Omega,
\end{align*}
with $p^S,p^X,p^P,p^\Phi$ supplied by~\eqref{eq:ac:pS}--\eqref{eq:ac:pPhi} and the five transversality conditions~\eqref{eq:ac:transversality} pinning down the initial values of the forward-looking variables. The instantaneous rates of change of the five stocks then follow from~\eqref{eq:ac:Phi-motion}, \eqref{eq:ac:reserves-motion}, \eqref{eq:ac:pollution-motion} and $\dot K=I$. Neither $\phi^Y$ nor $\phi^K$ is among the unknowns: the first is read off~\eqref{eq:ac:Omega} afterwards and the second off the state $\Phi$.

Under \textbf{(B1)} the system is \emph{nine} equations in nine unknowns. $\Omega$ leaves the allocation entirely --- \eqref{eq:ac:Omega} is no longer part of the simultaneous block but a post-processing step --- and \eqref{eq:ac:ledger-quadratic} is replaced by the closed form~\eqref{eq:ac:ledger-solved}. This is the single largest structural difference between the two settings, and the reason Section~\ref{sec:aq} adopts \textbf{(B1)}: it is the difference between solving for $W$ and reading it off.

It is worth recording which conditions are \emph{static} (they hold pointwise in $t$ given the stocks and the shadow prices) and which are \emph{intertemporal}: \eqref{eq:ac:foc-saving}, \eqref{eq:ac:foc-investment}, \eqref{eq:ac:foc-extraction}, \eqref{eq:ac:foc-exploration}, \eqref{eq:ac:foc-recycling}, \eqref{eq:ac:foc-treatment}, \eqref{eq:ac:pW-explicit}, \eqref{eq:ac:Omega} and~\eqref{eq:ac:ledger-quadratic} are static; \eqref{eq:ac:r}, \eqref{eq:ac:pS}--\eqref{eq:ac:pPhi} and~\eqref{eq:ac:keynes-ramsey} are intertemporal.

\smalltitle{Regimes.} The complementary-slackness structure of~\eqref{eq:ac:foc-recycling}, together with $c^T(0)=dc^T/d\varpi|_{\varpi=0}=0$, generates three regimes along the optimal path:
\begin{enumerate}[label=(\roman*),leftmargin=2.4em]
\item \emph{Linear stage:} $K^R=0$, $a=0$ and, if in addition $p^P=0$, condition~\eqref{eq:ac:foc-treatment} gives $dc^T/d\varpi=0$ and hence $\varpi=0$. Then $R^R=0$, $R=N$, and $p^W=c^c+p^P$ under \textbf{(B1)}, where in addition $W=N-\dot\Phi$; in general $W=N-\dot\Phi+\mathcal M$ and $p^W$ carries the denominator of~\eqref{eq:ac:pW-explicit}, which reduces to $1+\Omega\kappa\,\omega^Wc^c$ since $\varpi=0$.
\item \emph{Semi-linear stage:} $K^R=0$ but $\varpi>0$, sustained by $p^P>0$ through the term $p^P(1-d^W)$ in~\eqref{eq:ac:foc-treatment}. Waste is treated but not recycled.
\item \emph{Circular stage:} $K^R>0$ and $\varpi>0$, with~\eqref{eq:ac:foc-recycling-b} holding with equality.
\end{enumerate}
Regime~(ii) precedes~(iii) because $a''<0$ makes the marginal product of recycling capital low when the supply $\varpi W$ of recyclable material is small; and treatment is a precondition for recycling, so~(iii) cannot be reached without passing through a period with $\varpi>0$. The transition dates between regimes are outcomes of the model rather than parameters.

\subsection{Technical remarks}
\label{subsec:ac:remarks}

\begin{remark}[Why the materials balance is a ledger and not a constraint]\label{rem:ac:ledger}
Conservation of mass does not restrict what an economy can do; it records what the mass flows are given what the economy did. Nothing is gained by relaxing it, and no planner could. Attaching a multiplier to~\eqref{eq:ac:materials-balance} and treating it as a constraint would price a scarcity that the model already prices elsewhere --- reserves are finite ($p^S$), extraction is costly ($C^N_N$), and recycling recovers only $a(\cdot)$ of what is fed in.

The alternative reading, in which $\Omega$ and $\phi^Y$ are fixed parameters and~\eqref{eq:ac:materials-balance} restricts the allocation, is not available for a more elementary reason: with both fixed, \eqref{eq:ac:production-reduced} ties $R/Y$ to the constant $\Omega\omega^Y+\phi^Y$, contradicting the substitutability of capital for materials in~\eqref{eq:ac:output} and removing the mechanism the model is built to study. What \emph{is} restricted, and genuinely so, is that $\Omega\ge0$ in~\eqref{eq:ac:Omega}, i.e.\ $\phi^I(I+\delta K)\le N+R^R$: the economy cannot embody more mass in new capital than entered it. Equivalently $W\ge\delta\Phi$ in~\eqref{eq:ac:ledger}. This is generically slack --- it binds only in an implausibly fast accumulation --- but it is an inequality with a complementary-slackness multiplier, not an equality constraint with a multiplier that is nonzero by construction, and it should be monitored rather than imposed.
\end{remark}

\begin{remark}[The sign of $p^W$ is an outcome]\label{rem:ac:pW-sign}
Unlike $p^S,p^X,p^P$ and $p^\Phi$, the ledger price $p^W$ is not signed by the structure of the problem. From~\eqref{eq:ac:pW-explicit} --- whose denominator is positive in any sensible configuration, so the sign is carried by the numerator in the general case exactly as in \textbf{(B1)} --- $p^W<0$ whenever
\[
\varpi\,a(1-\varepsilon)\left[p^S+C^N_N+d^Wp^P\right]>c^c+c^T(\varpi)+p^P\left(1-\varpi+d^W\varpi\right),
\]
which is not exotic: it requires only that the recycling value of a marginal tonne exceed the cost of collecting, treating and disposing of it, and $p^S$ grows without bound as reserves are depleted. In that regime waste is a resource rather than a nuisance, $q>1$ by~\eqref{eq:ac:foc-investment}, and the Pigouvian tax on extraction implied by Section~\ref{sec:ad} turns into a subsidy --- extraction stocks the recycling loop, and that benefit is external to the extractor. This is a genuine feature of the accounting rather than a pathology, but it should be verified numerically and the complementarity solver must tolerate it.
\end{remark}

\begin{remark}[Bounds on the treatment share]\label{rem:ac:varpi}
Condition~\eqref{eq:ac:foc-treatment} is a necessary condition for an interior $\varpi\in(0,1)$ only. The assumption $dc^T/d\varpi|_{\varpi=0}=0$ in~\eqref{eq:ac:waste-cost} makes the lower bound non-binding: whenever the left side of~\eqref{eq:ac:foc-treatment} is strictly positive an interior solution exists, and when it is zero we get $\varpi=0$ as the solution of~\eqref{eq:ac:foc-treatment} itself. The upper bound is a genuine restriction: either $c^T$ must satisfy an Inada-type condition $\lim_{\varpi\to 1}dc^T/d\varpi=\infty$, or the constraint $\varpi\le 1$ must be imposed explicitly and~\eqref{eq:ac:foc-treatment} replaced by the corresponding complementary-slackness pair. The functional form adopted in the quantitative model does \emph{not} satisfy the Inada condition, so the bound is imposed there explicitly.
\end{remark}

\begin{remark}[The nature-attributable terms are what make $\Omega$ matter]\label{rem:ac:nature}
The economics behind \textbf{(B1)} is exact, and worth isolating from the algebra. Mass conservation pins waste down only when every gram of waste passed through $R$ on its way in. Overburden is a mass source outside the accounted flow, so it needs an intensity of its own, and that intensity has real allocative bite. This is why $\Omega$ is a pure reporting object under \textbf{(B1)} --- it multiplies nothing that the planner cares about --- and why it is not one in general.

What the restriction buys is visible in three places. The ledger becomes a division rather than a quadratic root, \eqref{eq:ac:ledger-solved} against~\eqref{eq:ac:ledger-quadratic}. The system loses an unknown and an equation, Section~\ref{subsec:ac:system}. And the relative intensities $\omega^i$ leave the allocation entirely, so their calibration becomes a reporting choice rather than a determinant of the transition date. What it costs is the physical content: overburden and gangue displaced by extraction, and any analogous byproduct of exploration, are no longer distinguished from waste generated by the final-good expenditure on those activities. Section~\ref{sec:aq} imposes the restriction and Remark~\ref{rem:aq:scope} records the consequences for the calibration exercise specifically.
\end{remark}

\begin{remark}[The dilution channel is an artifact, and what to do about it]\label{rem:ac:dilution}
The general accounting carries one channel that should be regarded with suspicion rather than defended. Because the nature-attributable flow is written $\mathcal M=\Omega\mathcal N$ --- overburden scaled by the economy-wide intensity level --- and because $\Omega=\left(R-\phi^IG\right)/\mathcal D$ is a \emph{residual ratio}, any activity that enlarges the base $\mathcal D$ mechanically lowers $\Omega$ and hence lowers overburden. This is the dilution term in~\eqref{eq:ac:OmegaN-derivative}, and it is what puts a wedge into the consumption condition~\eqref{eq:ac:foc-saving} and the factor $\left[1-\pi\omega^W\right]$ into~\eqref{eq:ac:foc-treatment}.

Read literally, it says that consuming more reduces mining overburden. That is not economics; it is an accounting identity being asked to carry a behavioral interpretation it was not built for. $\Omega$ is a \emph{measured} intensity --- waste per unit of final-good activity --- and the fact that a ratio falls when its denominator rises does not mean the numerator responded.

The clean fix is to scale the nature-attributable flow by an \emph{absolute} intensity rather than by $\Omega$: replace~\eqref{eq:ac:extraction} and~\eqref{eq:ac:exploration}'s second terms by $\bar\omega^{N,S}N$ and $\bar\omega^{D,S}D$ with $\bar\omega^{i,S}$ calibrated parameters in tonnes per tonne. Overburden per tonne of ore is exactly the sort of quantity material-flow accounts report directly, so this is if anything the better-identified specification. It also simplifies the theory considerably: $\mathcal M=\bar\omega^{N,S}N+\bar\omega^{D,S}D$ is then linear in the controls, the dilution term vanishes, $\pi\equiv0$, the ledger stays linear in $W$ rather than quadratic, and the only surviving general terms are the displacement ones --- $-p^W\bar\omega^{N,S}$ in~\eqref{eq:ac:foc-extraction} and $-p^W\bar\omega^{D,S}$ in~\eqref{eq:ac:foc-exploration}. Every dilution factor in Sections~\ref{subsec:ac:foc}--\ref{subsec:ac:costates} reduces to unity, and $B$ reduces to $\bar B+p^W\bar\omega^{N,S}$.

This document retains the $\Omega$-scaled form for continuity with the process-level accounting of Section~\ref{subsec:ac:accounting}, in which every intensity is relative and the level is derived. But the absolute specification is the one we would adopt if the general case were to be taken to data, and the choice should be settled before the nature-attributable terms are calibrated rather than after.
\end{remark}

\begin{remark}[Capital depreciation competes for recycling capacity]\label{rem:ac:depreciation}
Because depreciated capital enters $W$ on exactly the same footing as any other waste --- through the $\delta\Phi$ term in~\eqref{eq:ac:ledger}, equivalently through $-\dot\Phi$ --- replacing worn capital draws fresh raw material and the scrap it sheds competes for treatment share $\varpi$ and recycling capital $K^R$ on the same terms as waste from consumption or extraction. This is also what makes depreciation cost $q$ rather than $1$ in~\eqref{eq:ac:H-K}. An alternative, more restrictive assumption --- that depreciated capital material is costlessly and immediately re-embodied in new capital, bypassing $W$, $\varpi$ and $a(\cdot)$ entirely --- would leave capital replacement with no associated raw-material draw or waste at all, and would make $q\equiv1$. Whether that bypass is a reasonable simplification (capital scrap is disproportionately easy to recycle in practice) or a substantive mischaracterization is a modeling choice; this document adopts the process-level treatment, in which it is not assumed away.
\end{remark}

\begin{remark}[Why waste treatment needs only one channel]\label{rem:ac:units}
Extraction and exploration each split into a final-good-attributable component and a nature-attributable component, because tonnage extracted or discovered ($N$, $D$) and final-good expenditure on the activity ($C^N$, $C^D$) are physically distinct quantities that need not move together. Waste collection and treatment has no such second channel: there is no physical analogue of overburden for the act of treating waste, since the tonnage it processes, $W$, is by definition already accounted for elsewhere in the balance, and the only additional resource the activity draws on is the final good spent operating it, $C^W$. Hence $W^W=\Omega\omega^WC^W$ needs no further split.
\end{remark}

\begin{remark}[Interpretation of $\omega^{D,S}$]\label{rem:ac:dexploration}
The nature-attributable waste term for extraction, $\Omega\omega^{N,S}N$, has a direct physical reading: overburden and gangue displaced alongside the ore. The analogous term for exploration, $\Omega\omega^{D,S}D$, is less obviously a physical extraction from nature, since $D$ is new discoveries rather than a mined flow; it is retained for symmetry with $N$, but its calibration --- possibly at $\omega^{D,S}=0$ --- should be revisited once the two are compared against data.
\end{remark}

\begin{remark}[Instantaneous recirculation]\label{rem:ac:recirculation}
The circularity multiplier $\left(1-a\varpi\right)^{-1}$ in~\eqref{eq:ac:ledger-solved} has material passing through production, waste collection, treatment and recycling and back into production an unbounded number of times \emph{within the instant}. In continuous time this is innocuous --- it is the standard reading of a stock-flow identity with feedback --- but it becomes a substantive assumption once the model is discretized in Section~\ref{sec:aq}, where it means that a tonne extracted at the start of period $t$ can be recycled and re-used within period $t$. The natural alternative is a one-period lag on recycled input, $R^R_t=a_{t-1}\varpi_{t-1}W_{t-1}$, which removes the multiplier at the cost of an extra state. This document adopts the contemporaneous version throughout.
\end{remark}

\begin{remark}[One embodied-material stock, not two]\label{rem:ac:Phi-pools}
Equation~\eqref{eq:ac:pPhi} rests on the well-mixed assumption of Section~\ref{subsec:ac:embodied-stock}: $K^Y$ and $K^R$ are drawn from the same pool and share the average intensity $\phi^K(t)$, so that $\phi^K\delta K^Y+\phi^K\delta K^R=\delta\Phi$ and a single state $\Phi$ suffices. If that is relaxed --- recycling capital tracked as its own vintage pool with its own average intensity --- the co-state derivation would need to be redone with two embodied-material stocks and two liabilities, and the allocation of gross investment between the two pools would become a further margin. Nothing else in Sections~\ref{sec:ac}--\ref{sec:aq} depends on the assumption.
\end{remark}

\begin{remark}[Closing $\phi^I(t)$: two natural choices]\label{rem:ac:phiI}
$\phi^I(t)$ is left as a primitive process throughout Sections~\ref{sec:ac}--\ref{sec:ad}; nothing in the planner's problem, the decentralized economy, or the shadow-price system above depends on how it is closed. Two specifications are natural. \emph{(a) Exogenous:} $\phi^I(t)$ follows a prescribed calibrated path, independent of the allocation --- the direct generalization of treating $\phi^K$ as a constant, now allowing it to trend over time (e.g.\ a materials-saving or materials-worsening technology specific to capital-goods production). \emph{(b) Materials-linked:} $\phi^I(t)=\varsigma\,\Omega(t)$ for a constant $\varsigma$, tying the intensity of newly produced capital to the same economy-wide waste-intensity level $\Omega$ that governs every other final-good-denominated activity in Section~\ref{subsec:ac:accounting} --- capital-goods production is then simply $\varsigma$ times as material-intensive, per unit of final good, as activity in general; since $\Omega$ and $\phi^I$ share units, this is a dimensionally natural restriction, and $\varsigma=1$ is its most parsimonious case. Note that under rule~(b), and only under rule~(b), $\Omega$ re-enters the allocation even with $\omega^{N,S}=\omega^{D,S}=0$, since it then feeds~\eqref{eq:ac:ledger} through $\phi^I$. Both are pure closing rules: neither adds a control variable or a first-order condition to the system above, and both are stated explicitly in the quantitative implementation (Section~\ref{subsec:aq:phiI}).
\end{remark}

%%%%%%%%%%%%%%%%%%%%%%%%%%%%%%%%%%%%%%%%%%%%%%%%%%%%%%%
%% Alternative Part I, Section 2: Decentralized theory %%
%%%%%%%%%%%%%%%%%%%%%%%%%%%%%%%%%%%%%%%%%%%%%%%%%%%%%%%

\section{The Decentralized Model in Continuous Time}
\label{sec:ad}

This section states the market economy that corresponds to the planner model of Section~\ref{sec:ac}. Tastes and technologies are unchanged --- equations~\eqref{eq:ac:utility}--\eqref{eq:ac:waste-cost} carry over verbatim --- so this section describes only the institutional structure, the agents' problems, the resulting optimality conditions, and the policy instruments that implement the first best. The embodied-material stock $\Phi$ and its law of motion~\eqref{eq:ac:Phi-motion} also carry over unchanged: $\Phi$ is a real physical stock, not a planner construct, and the materials ledger~\eqref{eq:ac:ledger} determines actual total waste $W$ regardless of how policy is set.

One point of method is worth stating at the outset, because it differs from the usual treatment. The materials ledger closes \emph{identically} in the two economies: $\Omega$ and $\phi^Y$ are determined by~\eqref{eq:ac:Omega} in the market exactly as in the planner's problem, and $W$ by~\eqref{eq:ac:ledger-quadratic}. This is as it should be. Conservation of mass is physics, not an institution, and there is no sense in which a planner and a market face different versions of it. What differs between the two economies is only what the agents \emph{internalize}, and that is the entire subject of this section.

Like Section~\ref{sec:ac}, this section is stated under the general accounting; the restriction $\omega^{N,S}=\omega^{D,S}=0$ is displayed as \textbf{(B1)} wherever it simplifies a statement. It has a sharp consequence here: \textbf{(B1)} needs four instruments, and the general case needs seven. The three extra ones --- on output, on consumption and on exploration --- are all exactly zero under \textbf{(B1)}, which is why the baseline instrument set looks as parsimonious as it does. Throughout, $\kappa=\mathcal N/\mathcal D$ and $\pi=p^W\Omega\kappa$ as in Section~\ref{sec:ac}.

\subsection{Agents, markets and policy instruments}
\label{subsec:ad:setup}

The economy contains three decision units: a competitive waste industry, a representative household that owns and runs both the mining firm and the final goods firm, and a government. There is one market price in the model, the price $P^W\ge 0$ per unit of \emph{treated} waste sold by the waste industry to the final goods firm.

\smalltitle{The externalities to be corrected.} Private agents fail to internalize exactly two things: the pollution stock $P$, which the household takes parametrically, and the materials ledger~\eqref{eq:ac:ledger}, which no agent faces and which therefore has no private counterpart to $p^W$ or to the embodied-material liability $p^\Phi$. The reserve and discovery margins are already internalized, since the household owns the mine. Seven instruments suffice in general and four under \textbf{(B1)}, listed in Table~\ref{tab:ad:instruments}; all are allowed to be functions of time and are taken as given by private agents.

\begin{table}[H]
\centering
\caption{Policy instruments in the decentralized economy. The last column gives the instrument's value under \textbf{(B1)}; three of the seven vanish.}
\label{tab:ad:instruments}
\renewcommand{\arraystretch}{0.85}
\begin{tabular}{@{}lllll@{}}
\toprule
Instrument & Base & Levied on & Corrects & \textbf{(B1)} \\
\midrule
$\tau^N$        & $N$              & virgin materials extracted                & ledger, $p^W$ & $p^W$ \\
$s^R$           & $aW^R$           & waste actually recycled (subsidy)         & ledger and pollution & $d^Wp^P-p^W$ \\
$\tau^W$        & $(1-\varpi)W$    & untreated waste deposited in the environment & pollution, $p^P$ & $p^P(1-d^W)$ \\
$\tau^K$        & $I+\delta K$     & \emph{gross} capital formation            & ledger, $p^W$ and $p^\Phi$ & $-\phi^I\left(p^W-p^\Phi\right)$ \\
\midrule
$\tau^Y$        & $Y$              & output of final goods                     & dilution, $\pi$ & $0$ \\
$\tau^C$        & $C$              & consumption                               & dilution, $\pi$ & $0$ \\
$\tau^D$        & $D$              & new reserves discovered                   & displacement and dilution & $0$ \\
\midrule
$s^W$           & $W$              & waste collected (residual, from zero profit) & --- & --- \\
\bottomrule
\end{tabular}
\end{table}

$s^W$ is not a free instrument: it is whatever the competitive tender requires for the waste sector to break even, given the others and the market price $P^W$. Four features of the set are worth flagging now and are justified below.

First, the three instruments in the middle block exist \emph{only} because of the nature-attributable terms, and every one of them is proportional to $\pi$ or to $\Omega\omega^{i,S}$. Under \textbf{(B1)} they are identically zero, and the familiar four-instrument menu is what remains. It is worth being clear about which of the two readings is the surprising one: the absence of output and consumption taxes is a \emph{result} of \textbf{(B1)}, not a general property of materials-balance models.

Second, under \textbf{(B1)} there is no instrument levied on $Y$ or on $C$. Under a coarser accounting one expects ad valorem Pigouvian taxes on output and consumption, on the grounds that both generate waste. They do not, at the margin: by the ledger~\eqref{eq:ac:ledger}, $W$ then depends on $N$, $R^R$ and $\dot\Phi$ and on nothing else. The mass embodied in consumption goods was counted when it entered the economy as $R$; consuming it does not create it. In general $\tau^Y$ and $\tau^C$ reappear, but --- and this is the point --- \emph{not} as Pigouvian taxes on waste generation. They are dilution corrections, they carry the opposite sign, and Remark~\ref{rem:ad:no-tauY} argues they are an artifact of how $\mathcal M$ is scaled rather than a substantive result.

Third, there is no instrument levied on $C^N$, $C^D$ or $C^W$, and under \textbf{(B1)} $\tau^N$ carries no markup proportional to marginal extraction cost. Again this is the ledger at work: the final good spent on extraction effort carries embodied material, but that material was already drawn through $R$ and is already priced. In general $\tau^N$ does acquire a term in $C^N_N$, through dilution.

Fourth, $\tau^K$ is of a different character from all the others. It is not a Pigouvian tax correcting a waste- or pollution-generating externality; it is a price correction on the household's own saving decision, needed because capital is otherwise traded at a price that ignores the materials it embodies. It is levied on \emph{gross} capital formation $I+\delta K$, not on net investment, because that is the flow that draws material in the ledger and in~\eqref{eq:ac:Phi-motion}. And it turns out to be a \emph{subsidy}. Section~\ref{subsec:ad:foc} explains where it comes from, Section~\ref{subsec:ad:implementation} gives its Pigouvian level, and Remark~\ref{rem:ad:q} explains why no other instrument can substitute for it.

\subsubsection{The waste industry}
\label{subsubsec:ad:waste}

Collection and treatment of waste is delegated to specialized firms through a competitive public tender. Firms operate the constant-returns technology behind~\eqref{eq:ac:waste-cost}, are obliged to collect \emph{all} waste, may deposit untreated waste in the environment against the tax $\tau^W$ per unit, and receive the subsidy $s^W$ per unit collected. They take $W$ and $P^W$ as given and choose $\varpi$. Profit is
\begin{align}
\Pi^W &= \left[P^W\varpi+s^W-c^c-c^T(\varpi)-\tau^W(1-\varpi)\right]W,
\label{eq:ad:waste-profit}
\end{align}
with first-order condition
\begin{align}
\frac{dc^T}{d\varpi} &= P^W+\tau^W.
\label{eq:ad:foc-treatment}
\end{align}
The tender drives $\Pi^W$ to zero, which determines the subsidy required for the sector to break even:
\begin{align}
s^W &= c^c+c^T(\varpi)+\tau^W(1-\varpi)-P^W\varpi.
\label{eq:ad:zero-profit}
\end{align}
If $P^W$ becomes large enough, $s^W<0$ and the ``subsidy'' turns into a licence fee for the right to collect waste. Note that~\eqref{eq:ad:foc-treatment} and~\eqref{eq:ad:zero-profit} jointly imply the reduced form
\begin{align}
\varpi\frac{dc^T}{d\varpi} &= c^c+c^T(\varpi)+\tau^W-s^W.
\label{eq:ad:treatment-reduced}
\end{align}
Neither $\Omega$ nor any waste-intensity parameter enters the waste industry's problem: the firm sees only the physical technology $c^c+c^T(\varpi)$, exactly as in Section~\ref{sec:ac}'s primitives. The materials ledger is entirely an externality from the firm's point of view, internalized only through the instruments the government sets.

\subsubsection{The household, the mine and the final goods firm}
\label{subsubsec:ad:household}

The household has preferences~\eqref{eq:ac:utility}, owns the mine operating the cost functions~\eqref{eq:ac:extraction-cost}--\eqref{eq:ac:discovery-cost}, and owns the final goods firm operating~\eqref{eq:ac:output} and the recycling technology~\eqref{eq:ac:recycling}. The final goods firm buys the quantity $W^R$ of treated waste at price $P^W$ and recycles the fraction $a\!\left(K^R/W^R\right)$ of it; the firm's argument in the recycling function is $W^R$, whereas the planner's argument is $\varpi W$, and the two coincide in equilibrium by market clearing~\eqref{eq:ad:clearing} below. Since the waste industry earns zero profit, the household's net dividend is
\begin{align}
\Pi^H &= \left(1-\tau^Y\right)F\!\left(K-K^R,\ N+a\!\left(\tfrac{K^R}{W^R}\right)W^R,\ P\right)-C^N(N,S)-C^D(D,X)
\notag \\
&\phantom{{}={}}
+s^R a\!\left(\tfrac{K^R}{W^R}\right)W^R-\tau^NN-\tau^DD-P^WW^R.
\label{eq:ad:dividend}
\end{align}
Adding the lump-sum transfer $T$ and subtracting consumption at its tax-inclusive price, the household's budget constraint is
\begin{align}
Q\left(\dot K+\delta K\right) &= \Pi^H+T-\left(1+\tau^C\right)C,
&
Q &\equiv 1+\tau^K.
\label{eq:ad:budget}
\end{align}
Under \textbf{(B1)} the three instruments $\tau^Y$, $\tau^C$ and $\tau^D$ are zero and~\eqref{eq:ad:dividend}--\eqref{eq:ad:budget} reduce to their familiar forms.
$Q$ is the price at which the household acquires capital, in units of the final good. Because $\tau^K$ is levied on gross capital formation, replacement of depreciated capital is priced at $Q$ exactly as net accumulation is --- which is what the ledger requires, since the two draw material on identical terms.

\begin{problem}[Household]\label{prob:ad:household}
The household chooses $\left\{C,K^R,N,D,W^R\right\}$ to maximize~\eqref{eq:ac:utility} subject to~\eqref{eq:ad:budget}, $\dot S=D-N$ and $\dot X=D$, taking as given the paths of $P^W$, all policy instruments, the transfer $T$, and --- crucially --- the stock of pollution $P$. Initial stocks $K(0),S(0),X(0)$ are predetermined.
\end{problem}

The household internalizes neither the pollution externality nor the materials ledger. This is the entire source of the wedge between Problem~\ref{prob:ad:household} and Problem~\ref{prob:ac:planner}.

\smalltitle{Government budget.} The transfer is set to balance the budget period by period,
\begin{align}
T &= \tau^NN+\tau^DD+\tau^YY+\tau^CC+\tau^K\left(I+\delta K\right)+\tau^W(1-\varpi)W-s^WW-s^R aW^R.
\label{eq:ad:government-budget}
\end{align}
Substituting~\eqref{eq:ad:dividend} and~\eqref{eq:ad:government-budget} into~\eqref{eq:ad:budget}, every tax and subsidy term cancels. For $\tau^K$ this is immediate: the left side is $\left(1+\tau^K\right)\left(I+\delta K\right)$ and the transfer returns $\tau^K\left(I+\delta K\right)$, leaving $I+\delta K$. The same holds term by term for $\tau^Y$, $\tau^C$ and $\tau^D$, whose bases appear once in~\eqref{eq:ad:dividend} or~\eqref{eq:ad:budget} and once in~\eqref{eq:ad:government-budget} with the opposite sign. For the waste block, using~\eqref{eq:ad:zero-profit} and market clearing $W^R=\varpi W$,
\begin{align}
-P^W\varpi W+\tau^W(1-\varpi)W-s^WW &= -\left[c^c+c^T(\varpi)\right]W,
\label{eq:ad:waste-cancellation}
\end{align}
so the household budget collapses to the economy's resource constraint,
\begin{align}
\dot K &= F\!\left(K-K^R,R,P\right)-C-\delta K-C^N(N,S)-C^D(D,X)-\left[c^c+c^T(\varpi)\right]W,
\label{eq:ad:resource-constraint}
\end{align}
which is~\eqref{eq:ac:capital-constraint}. This is the check that the instrument set and the transfer rule are mutually consistent, and it holds with $\tau^K$ on the gross base exactly as it would with no capital instrument at all.

\subsection{Optimality conditions}
\label{subsec:ad:foc}

Because $\tau^K$ drives a wedge between the price of capital and the price of the final good, the household's problem carries two shadow prices where a market economy without $\tau^K$ needs only one. Let $\mu^K$ denote the true co-state on $K$ --- the household's marginal utility value of a unit of \emph{capital} --- and define
\begin{align}
\zeta^m &\equiv \frac{\mu^K}{Q},
\label{eq:ad:zetam}
\end{align}
the household's marginal utility value of a unit of the \emph{final good}, exactly mirroring the planner's $\zeta=\lambda^K/q$ from~\eqref{eq:ac:foc-investment}. Let $P^S>0$ and $P^X>0$ denote the private shadow prices of reserves and of accumulated discoveries, defined relative to $\zeta^m$ and measured in units of the final good. There is no private counterpart to $p^W$, to $p^\Phi$, or to $p^P$.

\smalltitle{The final goods firm.} Differentiating the recycling block of~\eqref{eq:ad:dividend} with respect to $K^R$ and $W^R$, and using $\partial\left[aW^R\right]/\partial W^R=a(1-\varepsilon)$. Writing $\widetilde F\equiv\left(1-\tau^Y\right)F$ for the after-tax marginal products,
\begin{subequations}\label{eq:ad:foc-recycling}
\begin{align}
K^R&=0
&&\text{if}\quad a'(0)\left[\widetilde F_R+s^R\right]\le \widetilde F_K,
&&\text{(linear stage)}
\label{eq:ad:foc-recycling-a}
\\
a'\!\left(\frac{K^R}{W^R}\right)\left[\widetilde F_R+s^R\right]&=\widetilde F_K
&&\text{if}\quad a'(0)\left[\widetilde F_R+s^R\right]> \widetilde F_K,
&&\text{(circular stage)}
\label{eq:ad:foc-recycling-b}
\end{align}
\end{subequations}
\begin{align}
a\!\left(\frac{K^R}{W^R}\right)(1-\varepsilon)\left[\widetilde F_R+s^R\right] &= P^W.
\label{eq:ad:foc-waste-use}
\end{align}

\smalltitle{Saving, extraction and exploration.}
\begin{align}
u'(C) &= \zeta^m\left(1+\tau^C\right),
&
\widetilde F_R-C^N_N-\tau^N &= P^S,
&
C^D_D+\tau^D+P^X &= P^S.
\label{eq:ad:foc-saving}
\end{align}
Under \textbf{(B1)}, $\tau^Y=\tau^C=\tau^D=0$, so $\widetilde F=F$ and these are the conditions of a market economy with no output, consumption or exploration instrument.

\smalltitle{Co-state conditions.} Differentiating~\eqref{eq:ad:budget} with respect to $K$, the household gives up $Q$ units of the final good per unit of capital and receives the after-tax marginal product $\widetilde F_K$ net of depreciation valued at $Q$, so
\begin{align}
\frac{\dot\mu^K}{\mu^K} &= \rho-r^m,
&
r^m &\equiv \frac{\widetilde F_K}{Q}-\delta,
&
\lim_{t\to\infty}\mu^K(t)K(t)&=0,
\label{eq:ad:muK}
\end{align}
which is~\eqref{eq:ac:r} with $Q$ in place of $q$. The goods rate of interest facing the household follows from $\zeta^m=\mu^K/Q$:
\begin{align}
\frac{\dot\zeta^m}{\zeta^m} &= \rho-\iota^m,
&
\iota^m &\equiv r^m+\frac{\dot Q}{Q},
\label{eq:ad:iotam}
\end{align}
mirroring~\eqref{eq:ac:iota}. $P^S$ and $P^X$, defined relative to $\zeta^m$ rather than $\mu^K$, therefore retain exactly their earlier form, discounted at $\iota^m$:
\begin{align}
P^S(t) &= \int_t^{\infty}\left(-C^N_S(z)\right)\exp\left(-\int_t^{z}\iota^m(u)\,du\right)dz,
&
P^X(t) &= \int_t^{\infty}C^D_X(z)\exp\left(-\int_t^{z}\iota^m(u)\,du\right)dz.
\label{eq:ad:PS-PX}
\end{align}
There is no private analogue of~\eqref{eq:ac:pP}: the household treats $P$ parametrically, so the pollution stock enters~\eqref{eq:ad:foc-recycling}--\eqref{eq:ad:foc-saving} only through the levels of $F_R$ and $F_K$ and never through a shadow price.

Every one of these conditions is the corresponding planner condition of Section~\ref{subsec:ac:foc} with the shadow prices $p^W$, $p^\Phi$ and $p^P$ deleted and the instruments put in their place. That is the whole content of the decentralization, and it is what makes the term-by-term matching in Section~\ref{subsec:ad:implementation} straightforward.

\subsection{Equilibrium}
\label{subsec:ad:equilibrium}

\begin{definition}[Competitive equilibrium]\label{def:ad:equilibrium}
Given initial stocks $K(0),\Phi(0),S(0),X(0),P(0)$ and paths of the four policy instruments, a competitive equilibrium is a path of quantities $\left\{C,K^R,N,D,W,W^R,\varpi\right\}$, a price $P^W\ge 0$, a subsidy $s^W$, and shadow prices $\left\{P^S,P^X,\mu^K\right\}$ (with $\zeta^m\equiv\mu^K/Q$ read off directly) such that
\begin{enumerate}[label=(\roman*),leftmargin=2.4em]
\item the waste industry optimizes and breaks even: \eqref{eq:ad:foc-treatment} and~\eqref{eq:ad:zero-profit};
\item the household and its firms optimize: \eqref{eq:ad:foc-recycling}--\eqref{eq:ad:PS-PX};
\item the market for treated waste clears, with the price non-negative and~\eqref{eq:ad:foc-waste-use} holding as an equality whenever the quantity traded is positive:
\begin{align}
\varpi W &= W^R\ge 0,
&
P^W &\ge 0,
&
P^W W^R &> 0 \ \Longrightarrow\ \eqref{eq:ad:foc-waste-use}\ \text{holds with equality};
\label{eq:ad:clearing}
\end{align}
\item total waste and the waste-intensity level satisfy the materials ledger and its definition, \eqref{eq:ac:ledger-quadratic} and~\eqref{eq:ac:Omega}, reducing under \textbf{(B1)} to $W=\left(N-\dot\Phi\right)/\left(1-a\varpi\right)$ from~\eqref{eq:ac:ledger-solved} with $\Omega$ dropping out;
\item the stocks evolve according to~\eqref{eq:ac:Phi-motion}, \eqref{eq:ac:reserves-motion}, \eqref{eq:ac:pollution-motion} and~\eqref{eq:ad:resource-constraint}, with $R^R=a\!\left(K^R/W^R\right)W^R$ in~\eqref{eq:ac:Xi}.
\end{enumerate}
The accounting object $\phi^Y$ is recovered from~\eqref{eq:ac:Omega} after the fact and enters no equilibrium condition; so does $\Omega$ under \textbf{(B1)}.
\end{definition}

\smalltitle{Counting.} Conditions~(i)--(iv) are nine equations in the nine unknowns
\begin{align*}
C,\quad K^R,\quad N,\quad D,\quad W,\quad W^R,\quad \varpi,\quad P^W,\quad \Omega,
\end{align*}
with $s^W$ read off~\eqref{eq:ad:zero-profit}, $P^S,P^X,\mu^K$ given by their own forward-looking relations~\eqref{eq:ad:muK}--\eqref{eq:ad:PS-PX}, and the transversality condition pinning down $\mu^K(0)$. $\Phi$ is predetermined at each instant and enters only through $\dot\Phi$ in~(iv); $\phi^K\equiv\Phi/K$ appears nowhere. Under \textbf{(B1)} the count falls to eight and eight: $\Omega$ leaves the system, exactly as in Section~\ref{subsec:ac:system}.

This is a genuinely square system under any policy configuration, not only the optimal one, and --- unlike a formulation in which the level $\Omega$ is a free calibrated parameter and the materials balance is imposed as an additional restriction --- nothing needs to be endogenized to make it so. Note the distinction: here $\Omega$ is an unknown \emph{determined by} an accounting identity, not a parameter whose value the modeller sets and against which the balance is then imposed as a restriction. The ledger supplies exactly one equation and determines exactly one quantity, $W$. Remark~\ref{rem:ad:cancellation} records why the cancellation that produces it is so clean.

\smalltitle{Regimes.} In the linear stage~\eqref{eq:ad:foc-recycling-a}, the demand for recyclable waste is zero. Then $W^R=0$, $P^W=0$, and all collected waste is deposited in the environment whether treated or not; treatment is positive only to the extent that $\tau^W>0$ makes it worthwhile through~\eqref{eq:ad:foc-treatment}. In the circular stage, \eqref{eq:ad:foc-recycling-b}, \eqref{eq:ad:foc-waste-use} and~\eqref{eq:ad:clearing} jointly determine $K^R$, $W^R$ and $P^W$. Under laissez-faire ($\tau^W=0$) the linear stage therefore has $\varpi=0$ exactly, so the model has a well-defined ``date of first waste treatment'' and a later ``date of first recycling.''

\subsection{Implementing the first best}
\label{subsec:ad:implementation}

\begin{proposition}[Pigouvian implementation]\label{prop:ad:implementation}
Let $p^W$, $p^S$, $p^X$, $p^P$, $p^\Phi$, $\kappa$, $\pi$ and $B$ from~\eqref{eq:ac:B-reduced} denote the shadow prices and accounting objects along the first-best path of Section~\ref{sec:ac}. Set
\begin{align}
\tau^Y &= -\pi\,\omega^Y,
&
\tau^C &= -\pi\,\omega^C,
&
\tau^D &= p^W\Omega\,\omega^{D,S}-\pi\,\omega^{D,Y}C^D_D,
\label{eq:ad:implement-general}
\\
\tau^N &= p^W\left(1+\kappa\right)+p^W\Omega\,\omega^{N,S}-\pi\,\omega^{N,Y}C^N_N,
&
s^R &= d^Wp^P-p^W\left(1+\kappa\right),
\notag
\\
\tau^W &= \frac{p^P\left(1-d^W\right)+\pi\,\omega^Wa\left(1-\varepsilon\right)B}{1-\pi\,\omega^W},
&
\tau^K &= -\phi^I\left[p^W\left(1+\kappa\right)-p^\Phi\right].
\notag
\end{align}
Then the competitive equilibrium of Definition~\ref{def:ad:equilibrium} coincides with the first-best allocation, and $Q=q$, $\mu^K=\lambda^K$, $\zeta^m=\zeta$, $P^S=p^S$, $P^X=p^X$, $r^m=r$, $\iota^m=\iota$ at every date. Under \textbf{(B1)}, where $\kappa=\pi=\Omega\omega^{N,S}=\Omega\omega^{D,S}=0$, this reduces to the four-instrument menu
\begin{align}
\tau^N &= p^W,
&
s^R &= d^Wp^P-p^W,
&
\tau^W &= p^P\left(1-d^W\right),
&
\tau^K &= -\phi^I\left(p^W-p^\Phi\right),
\label{eq:ad:implement}
\end{align}
with $\tau^Y=\tau^C=\tau^D=0$.
\end{proposition}

\smalltitle{Argument.} A term-by-term comparison; we record the six steps because each is a useful numerical consistency check.
\begin{enumerate}[label=\arabic*.,leftmargin=2.4em]
\item \emph{Output.} With $\tau^Y=-\pi\omega^Y$ we have $1-\tau^Y=1+\pi\omega^Y$, so $\widetilde F_K=F_K\left[1+\pi\omega^Y\right]=\widehat F_K$ and $\widetilde F_R=F_R\left[1+\pi\omega^Y\right]$, matching the factor that~\eqref{eq:ac:B} attaches to both marginal products. This step is what makes the remaining five work at all. It also explains the argument's structure under \textbf{(B1)}: there, a factor $\left(1-\tau^Y\right)$ would multiply $F_R$ and $F_K$ in~\eqref{eq:ad:foc-recycling} but only $F_R$ in the extraction condition, and the two cannot both be matched unless $\tau^Y=0$. In general the planner's own conditions carry exactly that asymmetric factor, which is why a nonzero $\tau^Y$ is now not merely admissible but required.
\item \emph{Recycling capital.} With $s^R=d^Wp^P-p^W\left(1+\kappa\right)$, the bracket in~\eqref{eq:ad:foc-recycling} becomes $\widetilde F_R+s^R=F_R\left[1+\pi\omega^Y\right]-p^W\left(1+\kappa\right)+d^Wp^P=B$, which is exactly~\eqref{eq:ac:B}, and the right-hand side is $\widetilde F_K=\widehat F_K$. Condition~\eqref{eq:ad:foc-recycling} therefore reproduces~\eqref{eq:ac:foc-recycling} including the switch point, so the market economy enters the circular stage on the same date as the planner.
\item \emph{The price of treated waste.} Substituting $s^R$ into~\eqref{eq:ad:foc-waste-use} gives $P^W=a(1-\varepsilon)B$, which by~\eqref{eq:ac:B-reduced} equals $a(1-\varepsilon)\left[\bar B+p^W\Upsilon\right]$, and under \textbf{(B1)} equals $a(1-\varepsilon)\left[p^S+C^N_N+d^Wp^P\right]>0$. Under \textbf{(B1)} the price of treated waste is therefore automatically strictly positive in the circular stage and no non-negativity restriction on it can bind; in general positivity requires $\bar B+p^W\Upsilon>0$, which should be checked rather than assumed.
\item \emph{Waste treatment.} Substituting $P^W$ from step~3 and $\tau^W$ into the waste industry's condition~\eqref{eq:ad:foc-treatment} gives $dc^T/d\varpi=\left[a(1-\varepsilon)B+p^P(1-d^W)\right]/\left(1-\pi\omega^W\right)$, which is the planner's~\eqref{eq:ac:foc-treatment} rearranged. Under \textbf{(B1)} the reading of $\tau^W$ is exact: it is the \emph{extra} pollution damage caused by depositing a tonne untreated rather than treated, $p^P$ against $d^Wp^P$. In general it also absorbs the dilution credit that treatment expenditure earns. The subsidy $s^W$ follows from~\eqref{eq:ad:zero-profit} and requires no separate matching, since it is pinned by zero profit rather than chosen.
\item \emph{Extraction and exploration.} With $\tau^N$ as given, the extraction condition in~\eqref{eq:ad:foc-saving} becomes $F_R\left[1+\pi\omega^Y\right]-C^N_N-\tau^N=P^S$, which is~\eqref{eq:ac:foc-extraction} given $P^S=p^S$ from step~6. Under \textbf{(B1)} exploration needs no instrument --- $C^D_D+P^X=P^S$ is already~\eqref{eq:ac:foc-exploration}, because a marginal discovery draws no material and generates no waste. In general it displaces mass ($\omega^{D,S}$) and dilutes the base ($\omega^{D,Y}$), and $\tau^D$ prices both.
\item \emph{Saving.} With $\tau^K=-\phi^I\left[p^W\left(1+\kappa\right)-p^\Phi\right]$ we have $Q=1+\tau^K=q$ exactly, hence $r^m=\widetilde F_K/Q=\widehat F_K/q=r$ and $\iota^m=\iota$ from~\eqref{eq:ad:muK}--\eqref{eq:ad:iotam} against~\eqref{eq:ac:r} and~\eqref{eq:ac:iota}. Note that this works only because step~1 put the dilution factor into $\widetilde F_K$ rather than into $Q$: had $\tau^Y$ been unavailable, matching $r^m$ to $r$ would have required $Q=qF_K/\widehat F_K$, whose time derivative then fails to match $\dot q/q$ in $\iota$, and no algebraic choice of $\tau^K$ would implement the optimum. Matching the transversality conditions and initial stocks gives $\mu^K\equiv\lambda^K$, hence $\zeta^m=\mu^K/Q=\lambda^K/q=\zeta$, hence $P^S=p^S$ and $P^X=p^X$ from~\eqref{eq:ad:PS-PX} against~\eqref{eq:ac:pS}--\eqref{eq:ac:pX}. Finally $u'(C)=\zeta^m\left(1+\tau^C\right)=\zeta\left[1-\pi\omega^C\right]$ is~\eqref{eq:ac:foc-saving}, and under \textbf{(B1)} reduces to $u'(C)=\zeta$ with no consumption tax required. Consistency with the government budget was verified in Section~\ref{subsubsec:ad:household}, and the identical allocations imply identical $\dot\Phi$, identical $W$ through the ledger, and identical pollution stocks. \hfill$\square$
\end{enumerate}

\smalltitle{Reading of the instrument set.} Take \textbf{(B1)} first, where the menu admits a one-line reading of each item.
\begin{itemize}[leftmargin=1.6em,itemsep=0.2em]
\item $\tau^N=p^W$ and $s^R=d^Wp^P-p^W$ are the \emph{same} instrument seen twice: a tax $p^W$ on every tonne of material entering production, $R=N+R^R$, with recycled material additionally credited $d^Wp^P$ for the pollution it avoids by not being deposited. Both bases carry the same materials charge because, by the ledger, a tonne is a tonne regardless of where it came from.
\item $\tau^W=p^P\left(1-d^W\right)$ is a pure pollution tax on the deposit margin.
\item $\tau^K=-\phi^I\left(p^W-p^\Phi\right)<0$ is a \emph{subsidy} to gross capital formation, equal to the social value of deferring waste by embodying it in a stock that decays at $\delta$ rather than releasing it now. By~\eqref{eq:ac:q-explicit} it is $-\phi^Ip^W\iota/(\iota+\delta)$ when $p^W$ is constant. It corrects no waste- or pollution-generating externality but a pure asset-pricing gap: capital would otherwise trade at one unit of the final good regardless of the material it embodies, \emph{over}stating its social cost.
\item $s^W$ is not Pigouvian at all. It exists because the waste sector is a regulated monopoly under tender, and it is pinned by zero profit rather than chosen.
\end{itemize}
A Pigouvian pollution tax alone is therefore not sufficient here, in contrast to models in which extraction of virgin material is the only source of pollution. But the reason is sharper than materials scarcity: it is that \emph{capital is a materials sink whose social value no market prices}.

The general menu preserves all four readings and modifies each in the same way. Every occurrence of the ledger price becomes $p^W\left(1+\kappa\right)$ --- a tonne drawn through production now drags $\kappa$ tonnes of overburden with it, so the materials charge rises proportionally, and it rises identically on $\tau^N$, $s^R$ and $\tau^K$, preserving the relations between them. On top of that, $\tau^N$ and $\tau^D$ acquire displacement terms in $\omega^{N,S}$ and $\omega^{D,S}$, which are genuinely Pigouvian: they price mass pulled out of the ground that no other instrument sees. Finally, $\tau^Y$, $\tau^C$ and the dilution parts of $\tau^N$, $\tau^D$ and $\tau^W$ are all proportional to $\pi$ and are of the opposite sign; these are the artifact discussed in Remark~\ref{rem:ad:no-tauY} and Remark~\ref{rem:ac:dilution}, and they would vanish identically under the absolute-intensity specification proposed there, leaving a general menu of five: the four of \textbf{(B1)}, with $p^W\left(1+\kappa\right)$ in place of $p^W$ and displacement terms added, plus $\tau^D$.

\subsection{Technical remarks}
\label{subsec:ad:remarks}

\begin{remark}[Why there is no output or consumption tax under \textbf{(B1)}]\label{rem:ad:no-tauY}
The absence of $\tau^Y$ and $\tau^C$ under \textbf{(B1)} is the most visibly counterintuitive feature of Proposition~\ref{prop:ad:implementation} and deserves a direct statement. It is not an assumption that consumption is clean. It is the ledger: by~\eqref{eq:ac:ledger}, total waste is then $R-\dot\Phi$, so the derivative of $W$ with respect to $C$, holding $N$, $R^R$ and $\dot\Phi$ fixed, is zero. Mass is conserved, and consuming a good does not bring new mass into existence --- the mass was priced when it was extracted or recycled, through $\tau^N$ and $s^R$.

What consumption \emph{does} affect is the composition of $Y$ between consumption and investment, and that margin is priced, by $\tau^K$: shifting the bundle toward investment stores mass, and shifting it toward consumption releases mass. Taxing $C$ and subsidising $I+\delta K$ would be double counting the same margin. The empirical content of the $\omega^i$'s is not lost either; they are what determine $\Omega$ and $\phi^Y$ in~\eqref{eq:ac:Omega}, and hence what the model says about material intensities in the data. They simply do not enter the allocation, unless the nature-attributable terms are retained or $\phi^I$ is closed by rule~(b) of Remark~\ref{rem:ac:phiI}.

That the general case restores $\tau^Y$ and $\tau^C$ does not overturn any of this, and it would be a mistake to read~\eqref{eq:ad:implement-general} as vindicating the intuition that output and consumption should be taxed because they generate waste. The restored instruments are \emph{subsidies}, not taxes, whenever $p^W>0$; they are proportional to $\pi$, hence to the dilution channel and to nothing else; and they are zero when extraction and exploration displace no mass, however waste-intensive output and consumption may be. Their content is that enlarging the denominator of a ratio lowers the ratio --- Remark~\ref{rem:ac:dilution} --- and under the absolute-intensity specification proposed there they vanish identically while the displacement terms in $\tau^N$ and $\tau^D$ survive. The correct summary is therefore unchanged: \emph{no instrument on output or consumption corrects a waste-generation externality in this model, under either accounting.}
\end{remark}

\begin{remark}[Why the ledger collapses so cleanly]\label{rem:ad:cancellation}
The collapse of~\eqref{eq:ac:waste-generation} and~\eqref{eq:ac:materials-balance} to the ledger~\eqref{eq:ac:ledger} happens because the two share the \emph{same} weighted sum $\mathcal D$ on their right-hand sides, differing only in how capital is treated: the materials balance charges gross investment $G=I+\delta K$ at the current intensity $\phi^I$, while the waste-generation identity credits depreciation $\delta\Phi$ at whatever intensity the stock carries. Substituting the first into the second cancels $\Omega\mathcal D$ and with it every $\omega^i$ appearing in $\mathcal D$, and what is left of the capital terms, $-\phi^IG+\delta\Phi$, is exactly $-\dot\Phi$ by~\eqref{eq:ac:Phi-motion}. The apparent complexity of the vintage structure --- $\phi^I$ versus $\phi^K$ --- disappears entirely into a single state variable's rate of change.

What does \emph{not} cancel is the second, independent residual $\Omega\mathcal N$, because $\mathcal N$ appears in the waste-generation identity and not in the materials balance --- that mass was never backed by material drawn through $R$, so there is nothing on the other side for it to cancel against. Under \textbf{(B1)} it is absent and the ledger is a division; in general it survives and the ledger is a quadratic, \eqref{eq:ac:ledger-quadratic}. Everything that distinguishes the two settings, in all three sections of this document, traces back to this one term.
\end{remark}

\begin{remark}[Why the capital-pricing gap needed a price, not a quantity, instrument]\label{rem:ad:q}
The household's budget constraint prices capital at one unit of the final good regardless of $\phi^I$, $p^W$ or $p^\Phi$, so $\mu^K$ never acquires the wedge $q=1-\phi^I\left(p^W-p^\Phi\right)$ that separates $\zeta$ from $\lambda^K$ in~\eqref{eq:ac:foc-investment}. $\tau^K$ closes this by construction: it is levied on the margin where the wedge lives --- gross capital formation, the flow that actually draws material --- at exactly the rate the planner's own first-order condition for $I$ implies, so it reproduces $q$ exactly rather than approximately or on average.

Two features are worth separating. First, the \emph{base} matters. Levying $\tau^K$ on net investment $I$ alone would leave replacement of depreciated capital priced at one, whereas the ledger charges it at $\phi^I$ like any other capital formation; the resulting mismatch is exactly the term $\delta\left(q-1\right)$ that distinguishes a correct return $r=F_K/q-\delta$ from the naive $F_K-\delta$. Second, no quantity restriction could substitute. Forcing $K^R$, $N$, or any other control to a particular path cannot correct a distortion that lives in a \emph{price} the household faces, because the household would simply re-optimize every other margin against the same wrong relative price of capital and goods.
\end{remark}

\begin{remark}[Which shadow prices the market model needs]\label{rem:ad:pP}
None of $p^P$, $p^W$ or $p^\Phi$ appears anywhere in conditions~(i)--(v) of Definition~\ref{def:ad:equilibrium}. All three are purely planner-side objects with no private counterpart, and they enter the market economy \emph{only} through the policy rules~\eqref{eq:ad:implement-general}. Consequently: under exogenous policy the market model can be solved without ever computing any of them; under optimal policy all three must be computed alongside the market allocation. Section~\ref{subsec:aq:optimal-policy} writes this out in discrete time.

$\Omega$ is different, and the difference is exactly \textbf{(B1)}. Under \textbf{(B1)} it is like the three planner prices --- absent from the equilibrium conditions, needed only to report the instruments in interpretable units. In general it is an equilibrium unknown in its own right, present in condition~(iv) under \emph{any} policy including laissez-faire, and therefore not something the exogenous-policy solver can defer.
\end{remark}

\begin{remark}[Two arguments of the recycling function]\label{rem:ad:argument}
The planner writes the recycling rate as $a\!\left(K^R/(\varpi W)\right)$, the firm as $a\!\left(K^R/W^R\right)$. These agree in equilibrium by~\eqref{eq:ad:clearing}, but the derivatives do not agree \emph{off} equilibrium, which is exactly why the planner's~\eqref{eq:ac:foc-treatment} contains a term the firm's condition~\eqref{eq:ad:foc-treatment} does not: the planner internalizes that treating more waste raises the supply of recyclable material, while the firm sees only $P^W+\tau^W$. The market price $P^W$ is what transmits this margin, and $\tau^W$ is what corrects the residual pollution wedge. For the numerical implementation the practical implication is that the recycling block should be written in terms of the capital intensity
\begin{align}
x &\equiv \frac{K^R}{W^R} = \frac{K^R}{\varpi W}
\label{eq:ad:x}
\end{align}
as a single unknown, rather than as a ratio of two quantities that are both zero in the linear stage.
\end{remark}

\begin{remark}[Instrument non-uniqueness]\label{rem:ad:normalization}
Proposition~\ref{prop:ad:implementation} states one member of a family. Adding a tax $\tau^{W^R}$ on purchases of treated waste $W^R$ changes~\eqref{eq:ad:foc-waste-use} to $a(1-\varepsilon)\left[F_R+s^R\right]=P^W+\tau^{W^R}$ and lowers the equilibrium $P^W$ one for one, so only the difference $\tau^W-\tau^{W^R}=p^P\left(1-d^W\right)$ is pinned down, not the two separately. Setting $\tau^{W^R}=0$, as above, is the normalization that makes $\tau^W$ readable as a pure pollution tax and keeps $P^W$ at its shadow value $a(1-\varepsilon)B$; any other split implements the same allocation with a different, and less interpretable, decomposition between the price and the tax. The same is true of transfers between $s^W$ and $\tau^W$, which~\eqref{eq:ad:zero-profit} absorbs. Only the instruments' \emph{net} effect on each margin is an economic object.

The general case adds a second, larger instance of the same non-uniqueness, and it is worth naming because it affects how the results should be reported. The dilution corrections enter every margin through the common factor $\pi$, and they can be levied on the activity or subtracted from its counterpart --- $\tau^Y$ against a uniform scaling of $\tau^N$, $s^R$ and $\tau^K$, for instance --- with no change in the allocation. The normalization chosen in Proposition~\ref{prop:ad:implementation} puts the whole of the output-side dilution into $\tau^Y$, which keeps $Q=q$ and hence keeps every shadow-price identity of Section~\ref{sec:ac} intact in the market economy. That is the right choice for verification, since it maximizes the number of quantities that must agree across the two solutions; it is not necessarily the right choice for presenting tax rates to a reader.
\end{remark}

%%%%%%%%%%%%%%%%%%%%%%%%%%%%%%%%%%%%%%%%%%%%%%%%%%%%%%%
%%    Alternative Part I, Section 3: Quantitative      %%
%%%%%%%%%%%%%%%%%%%%%%%%%%%%%%%%%%%%%%%%%%%%%%%%%%%%%%%

\section{The Quantitative Model in Discrete Time}
\label{sec:aq}

This section states the discrete-time version of the decentralized model of Section~\ref{sec:ad} with explicit functional forms, and reduces it to the system that is actually solved. It also presents a number of restricted variants that are useful for setting up, testing and debugging the numerical implementation. The numerical algorithms themselves are described in Part~\ref{part:num}.

Two features drive most of the results: \emph{technical progress}, which enters every technology and cost function at a sector-specific rate, and a \emph{minimum materials requirement} $\bar R$ in the production function, which bounds the substitutability of man-made for natural capital. A third feature is specific to this document: because the materials balance is a ledger rather than a constraint (Section~\ref{subsec:ac:balances}), the waste-intensity \emph{level} $\Omega_t$ is neither a calibrated parameter nor an equilibrium unknown but a derived reporting object, and total waste $W_t$ is pinned down by conservation of mass alone.

\smalltitle{One restriction, imposed here and nowhere else.} Sections~\ref{sec:ac} and~\ref{sec:ad} are stated under the general accounting. This section imposes
\begin{align}
\omega^{N,S} &= \omega^{D,S} = 0
&&\textbf{(B1)}
\label{eq:aq:B1}
\end{align}
throughout, and it is the \emph{only} departure from the general theory made in the quantitative model. Section~\ref{subsec:aq:B1} states precisely what \textbf{(B1)} buys and Remark~\ref{rem:aq:scope} what it costs. Everything else in this section is the general model: no other simplification is made silently, and every remaining modelling choice is a switch catalogued in Section~\ref{sec:baseline} with its alternative stated.

Note on notation: $\phi$ without a superscript is the curvature of the disutility of pollution in~\eqref{eq:aq:utility}; $\phi^I_t$, $\phi^K_t$ and $\phi^Y_t$ are the material intensities of Section~\ref{sec:ac} and are unrelated to it.

\subsection{Timing and states}
\label{subsec:aq:timing}

At the beginning of period $t$ the five stocks $\left(K_t,\Phi_t,S_t,X_t,P_t\right)$ are predetermined. Within the period the flows $\left(C_t,Y_t,N_t,D_t,W_t,W^R_t,\varpi_t,K^R_t\right)$ are chosen. Writing gross capital formation as
\begin{align}
G_t &\equiv I_t+\delta K_t = K_{t+1}-\left(1-\delta\right)K_t,
&
I_t &\equiv K_{t+1}-K_t,
\label{eq:aq:gross-investment}
\end{align}
the stocks are carried forward according to
\begin{align}
G_t &= Y_t-C_t-C^N_t-C^D_t-\left[c^c_t+c^T_t(\varpi_t)\right]W_t,
\label{eq:aq:K-motion}
\\
\Phi_{t+1} &= (1-\delta)\Phi_t+\phi^I_tG_t,
\qquad\qquad
\Delta\Phi_t \equiv \Phi_{t+1}-\Phi_t = \phi^I_tG_t-\delta\Phi_t,
\label{eq:aq:Phi-motion-discrete}
\\
S_{t+1} &= S_t+D_t-N_t,
\qquad\qquad
X_{t+1} = X_t+D_t,
\label{eq:aq:SX-motion}
\\
P_{t+1} &= \left[1-\theta(P_t)\right]P_t+\Xi_t,
\qquad
\Xi_t = W_t\left(1-\varpi_t+d^W\varpi_t\right)-d^WR^R_t.
\label{eq:aq:P-motion}
\end{align}
The average intensity of the capital stock is the derived ratio $\phi^K_t\equiv\Phi_t/K_t$; \eqref{eq:aq:Phi-motion-discrete} is the discrete-time counterpart of~\eqref{eq:ac:Phi-motion}, and $\phi^I_t$ is specified in Section~\ref{subsec:aq:phiI}. Since $P_t$ is predetermined, the productivity feedback from pollution operates with a one-period lag; this is what keeps the within-period system block-recursive in $P$.

\subsection{Functional forms}
\label{subsec:aq:functional-forms}

\smalltitle{Preferences.} Utility is CRRA in consumption with a separable, convex disutility of pollution whose weight grows over time as the marginal willingness to pay for environmental quality rises with income:
\begin{align}
U_0 &= \sum_{t=0}^{\infty}\beta^t\left[\frac{C_t^{1-\sigma}}{1-\sigma}-\frac{\psi_t}{1+\phi}P_t^{1+\phi}\right],
&
\psi_t &= \bar\psi\left(1+g^\psi\right)^t,
&
\beta&=\frac{1}{1+\rho}\in(0,1),
\quad \sigma,\phi,\bar\psi>0,\ g^\psi\ge 0.
\label{eq:aq:utility}
\end{align}

\smalltitle{Final goods and recycling.}
\begin{align}
Y_t &= A_t\left(K^Y_t\right)^{\alpha}\left(R_t-\bar R\right)^{1-\alpha},
&&
0<\alpha<1,\quad R_t\ge\bar R,\quad \bar R\ge 0,
\label{eq:aq:production}
\\
A_t &= \bar A\left(1+g^A\right)^tP_t^{-\gamma},
&&
\gamma>0,
\label{eq:aq:productivity}
\\
K^Y_t &= K_t-K^R_t,
&&
R_t = N_t+a_tW^R_t,
\label{eq:aq:definitions}
\\
a_t &= a^m_t\frac{x_t}{1+x_t},
&&
x_t\equiv\frac{K^R_t}{W^R_t},
\label{eq:aq:recycling-rate}
\\
a^m_t &= 1-\frac{b}{\left(1+g^R\right)^t},
&&
0<b<1,\quad g^R\ge 0.
\label{eq:aq:recycling-frontier}
\end{align}
$\bar R$ is the minimum materials input required by the laws of physics; $\bar R=0$ returns the standard Cobb--Douglas case. The maximum feasible recycling rate $a^m_t$ rises towards but never reaches unity, which is what keeps the circularity multiplier $\left(1-a_t\varpi_t\right)^{-1}$ of~\eqref{eq:ac:ledger-solved} finite.

\smalltitle{Extraction, exploration and waste costs.}
\begin{align}
C^N_t &= \frac{c^N_t}{1+\varepsilon^S}N_tS_t^{-\left(1+\varepsilon^S\right)},
&
c^N_t &= \frac{c^N}{\left(1+g^N\right)^t},
&
\varepsilon^S&>0,\ g^N\ge 0,
\label{eq:aq:extraction-cost}
\\
C^D_t &= \frac{c^D_t}{1+\varepsilon^X}D_tX_t^{1+\varepsilon^X},
&
c^D_t &= \frac{c^D}{\left(1+g^D\right)^t},
&
\varepsilon^X&>0,\ g^D\ge 0,
\label{eq:aq:discovery-cost}
\\
c^c_t &= \frac{\bar c^c}{\left(1+g^c\right)^t},
&
c^T_t(\varpi) &= \frac{\bar c^T}{\left(\varepsilon^T+1\right)\left(1+g^T\right)^t}\varpi^{\varepsilon^T+1},
&
\varepsilon^T&>0,\ g^c,g^T\ge 0.
\label{eq:aq:waste-cost}
\end{align}

\smalltitle{Waste intensities and absorption.} The relative intensities are calibrated technological parameters; the level $\Omega_t$ is derived in Section~\ref{subsec:aq:materials} and is deliberately absent from this block:
\begin{align}
\omega^i_t &= \frac{\omega^i}{\left(1+g^i_\omega\right)^t},\quad i\in\{Y,N,D,C,W\},
&
\theta(P_t) &= \bar\theta P_t^{-\varepsilon^P},
\quad 0<\varepsilon^P<1.
\label{eq:aq:intensities}
\end{align}
$\omega^N_t,\omega^D_t$ here are what Section~\ref{sec:ac} calls $\omega^{N,Y}_t,\omega^{D,Y}_t$; the nature-attributable components are set to zero throughout this document.

\smalltitle{Derivatives.} Collecting the derivatives that appear in the equilibrium conditions,
\begin{align}
F_{K,t} &= \alpha\frac{Y_t}{K^Y_t},
&
F_{R,t} &= (1-\alpha)\frac{Y_t}{R_t-\bar R},
&
F_{P,t} &= -\gamma\frac{Y_t}{P_t},
\label{eq:aq:F-derivatives}
\\
C^N_{N,t} &= \frac{c^N_t}{1+\varepsilon^S}S_t^{-\left(1+\varepsilon^S\right)},
&
C^N_{S,t} &= -c^N_tN_tS_t^{-\left(2+\varepsilon^S\right)},
&
&
\label{eq:aq:CN-derivatives}
\\
C^D_{D,t} &= \frac{c^D_t}{1+\varepsilon^X}X_t^{1+\varepsilon^X},
&
C^D_{X,t} &= c^D_tD_tX_t^{\varepsilon^X},
&
\frac{dc^T_t}{d\varpi} &= \frac{\bar c^T}{\left(1+g^T\right)^t}\varpi_t^{\varepsilon^T},
\label{eq:aq:CD-derivatives}
\\
\widehat\theta_t &\equiv \frac{d\left[\theta(P_t)P_t\right]}{dP_t}=\bar\theta P_t^{-\varepsilon^P}\left(1-\varepsilon^P\right)>0.
\label{eq:aq:theta-hat}
\end{align}
$C^N_t$ and $C^D_t$ are \emph{linear} in the flows $N_t$ and $D_t$, so $C^N_{NN}=C^D_{DD}=0$; this is the boundary case of the assumptions~\eqref{eq:ac:extraction-cost}--\eqref{eq:ac:discovery-cost} and has a consequence for the exploration margin recorded in Remark~\ref{rem:aq:linear-costs}.

\subsection{Properties of the recycling technology}
\label{subsec:aq:recycling-properties}

Writing the recycling block in terms of the capital intensity $x_t$ rather than in terms of $K^R_t$ and $W^R_t$ separately (Remark~\ref{rem:ad:argument}) makes the whole block closed-form. From~\eqref{eq:aq:recycling-rate},
\begin{align}
\varepsilon_t &\equiv a'_t\frac{x_t}{a_t}=\frac{1}{1+x_t}\in(0,1),
&
a_t &= a^m_t\left(1-\varepsilon_t\right),
&
a'_t &= a^m_t\varepsilon_t^2=\frac{a^m_t}{\left(1+x_t\right)^2},
&
a_t\left(1-\varepsilon_t\right) &= a^m_t\frac{x_t^2}{\left(1+x_t\right)^2}.
\label{eq:aq:recycling-algebra}
\end{align}
In particular $a'_t(0)=a^m_t$, and the ratio $a_t(1-\varepsilon_t)/a'_t=x_t^2$ exactly --- the identity that collapses the recycling block below. All the assumptions in~\eqref{eq:ac:recycling} hold: $a_t(0)=0$, $a'_t>0$, $a''_t<0$, $\varepsilon_t<1$, and $a_t\to a^m_t<1$ as $x_t\to\infty$.

\subsection{The baseline restriction \textbf{(B1)}}
\label{subsec:aq:B1}

Before stating the equilibrium system it is worth setting out exactly what~\eqref{eq:aq:B1} does, since every simplification in the rest of this section descends from it. Under \textbf{(B1)} the nature-attributable base $\mathcal N_t$ is identically zero, hence so are $\kappa_t=\mathcal N_t/\mathcal D_t$ and the dilution credit $\pi_t=p^W_t\Omega_t\kappa_t$, and with them:

\begin{enumerate}[label=(\arabic*),leftmargin=2.4em,itemsep=0.25em]
\item \emph{The ledger becomes a division.} The quadratic~\eqref{eq:ac:ledger-quadratic} factors, and $W_t$ is given in closed form by~\eqref{eq:aq:waste} below. This is the substantive gain: $W_t$ is read off rather than solved for.
\item \emph{$\Omega_t$ leaves the equilibrium system.} It becomes a post-processing step, computed from~\eqref{eq:aq:Omega} once the allocation is known --- except under closing rule~(b) of Section~\ref{subsec:aq:phiI}, which reintroduces it by a different route. The period's unknown count falls by one.
\item \emph{The relative intensities $\omega^i$ leave the allocation.} They determine the reported $\Omega_t$ and $\phi^Y_t$ and nothing else, again except under rule~(b). This is what makes their calibration a reporting choice; see Section~\ref{subsec:aq:calibration}.
\item \emph{Three instruments vanish.} $\tau^Y_t=\tau^C_t=\tau^D_t=0$ in Proposition~\ref{prop:ad:implementation}, and the remaining four take the simple forms~\eqref{eq:ad:implement}. The seven-instrument general menu collapses to four.
\item \emph{Every dilution factor becomes unity.} $B_t$ reduces to~\eqref{eq:ac:B-reduced-B1}, $p^W_t$ to~\eqref{eq:ac:pW-explicit-B1}, $\widehat F_{K,t}$ to $F_{K,t}$, and the marginal products in the recycling and extraction blocks lose their $\left[1+\pi_t\omega^Y\right]$ factors. This is what preserves the closed forms of Section~\ref{subsec:aq:system}.
\end{enumerate}

Point~(1) is the one that matters for the numerical strategy of Part~\ref{part:num}; points~(3) and~(5) are the ones that matter for how the results should be interpreted. Nothing in Part~\ref{part:num} depends on \textbf{(B1)} beyond point~(1).

\subsection{Materials, waste, and the derived $\Omega_t$}
\label{subsec:aq:materials}

Under \textbf{(B1)} the materials ledger~\eqref{eq:ac:ledger} in discrete time reads $W_t=N_t+R^R_t-\Delta\Phi_t$, and since $R^R_t=a_t\varpi_tW_t$ it solves explicitly:
\begin{align}
\boxed{\;W_t = \frac{N_t-\Delta\Phi_t}{1-a_t\varpi_t}
= \frac{N_t-\phi^I_tG_t+\delta\Phi_t}{1-a_t\varpi_t}\;}
\label{eq:aq:waste}
\end{align}
The numerator is the direct implication of conservation of mass: material extracted this period, less what stayed behind in the capital stock net of what the stock released through depreciation. The denominator is the circularity multiplier. Both objects in the numerator are unambiguous, and neither $\Omega_t$ nor any $\omega^i_t$ appears.

The level $\Omega_t$ and the material intensity $\phi^Y_t$ of the final good are recovered afterwards from~\eqref{eq:ac:Omega}:
\begin{align}
\Omega_t &= \frac{N_t+R^R_t-\phi^I_tG_t}{\mathcal D_t},
&
\mathcal D_t &\equiv \omega^Y_tY_t+\omega^N_tC^N_t+\omega^D_tC^D_t+\omega^C_tC_t+\omega^W_tC^W_t,
&
C^W_t &\equiv \left[c^c_t+c^T_t(\varpi_t)\right]W_t.
\label{eq:aq:Omega}
\end{align}
$\Omega_t$ is not among the period's unknowns (Section~\ref{subsec:aq:system}) and plays no role in determining the allocation. It matters for two things only: translating the allocation into the tax rates of Section~\ref{subsec:aq:optimal-policy} in interpretable units, and reporting the implied path of the economy's waste intensity as a check on the calibration. Under closing rule~(b) of Section~\ref{subsec:aq:phiI} it acquires a third role, feeding back into $\phi^I_t$. $\phi^K_t=\Phi_t/K_t$ is likewise a reporting object: it is predetermined, but it enters~\eqref{eq:aq:waste} only through $\delta\Phi_t$ and feeds back into none of the period's other equilibrium conditions.

\subsection{Closing $\phi^I_t$}
\label{subsec:aq:phiI}

$\phi^I_t$ has been left unspecified up to this point, exactly as in Sections~\ref{sec:ac}--\ref{sec:ad} (Remark~\ref{rem:ac:phiI}): neither closing rule below adds a control variable or a first-order condition.

\smalltitle{Neither rule is the baseline.} Unlike \textbf{(B1)}, which is a restriction of the general model and is imposed here once and for all, the two rules below are a genuine \emph{fork}: neither nests the other, neither is more general, and this document takes no position on which is right. Both are carried through the quantitative exercise and results are reported under both --- Section~\ref{subsec:aq:scenarios} makes the comparison an experiment in its own right. The distinction matters because the two rules disagree about something substantive: under~(a) the relative intensities $\omega^i$ do not touch the allocation, and under~(b) they do. Reporting only one rule would therefore not be a simplification but a hidden answer to an open question.

\smalltitle{(a) Exogenous.} $\phi^I_t$ follows a prescribed path, independent of the allocation, in the same family as the other technology parameters:
\begin{align}
\phi^I_t &= \frac{\bar\phi^I}{\left(1+g^\phi\right)^t},
&
g^\phi \gtrless 0,
\label{eq:aq:phiI-exogenous}
\end{align}
with $g^\phi>0$ a materials-saving trend in capital-goods production and $g^\phi<0$ the reverse. $g^\phi=0$ recovers a constant flow intensity, though not a constant $\phi^K_t$: the stock still converges toward it via~\eqref{eq:aq:Phi-motion-discrete} unless $\Phi_0=\bar\phi^IK_0$ exactly. Under this rule $\Omega_t$ is a pure reporting device, computed after the fact.

\smalltitle{(b) Materials-linked.} $\phi^I_t=\varsigma\,\Omega_t$ for a constant $\varsigma$ ties new capital's intensity to the same economy-wide level $\Omega_t$ that governs every other final-good-denominated activity. Substituting into~\eqref{eq:aq:Omega} and solving --- the self-reference is purely contemporaneous, so this remains closed-form:
\begin{align}
\Omega_t &= \frac{N_t+R^R_t}{\mathcal D_t+\varsigma G_t},
&
\phi^I_t &= \varsigma\,\Omega_t.
\label{eq:aq:phiI-linked}
\end{align}
Unlike under rule~(a), $\Omega_t$ here is not a pure reporting device: it feeds into~\eqref{eq:aq:phiI-linked} first, which then feeds into~\eqref{eq:aq:Phi-motion-discrete} and~\eqref{eq:aq:waste}, so the relative intensities $\omega^i$ become allocation-relevant. It remains fully closed-form given the period's other flows, so it adds no residual unknown to the count in Section~\ref{subsec:aq:system}, but it does add a dependence that rule~(a) does not have.

Either rule leaves $\Phi_t$ predetermined at the start of period $t$, so the recycling block, the extraction and exploration blocks, and the Euler equation are unaffected by the choice between them; only~\eqref{eq:aq:Phi-motion-discrete}, \eqref{eq:aq:waste}, and the level of $\tau^K_t$ depend on $\phi^I_t$ directly.

\subsection{The equilibrium system}
\label{subsec:aq:system}

We state the market-economy conditions of Definition~\ref{def:ad:equilibrium} in discrete time. Write $Q_t\equiv1+\tau^K_t$ for the price of capital in units of the final good.

\smalltitle{The goods rate of interest.} A unit of the final good at $t$ buys $1/Q_t$ units of capital, which yields the marginal product $F_{K,t+1}$ at $t+1$ and leaves $\left(1-\delta\right)$ units worth $Q_{t+1}$ each. Hence the discrete-time counterpart of~\eqref{eq:ad:iotam} is
\begin{align}
\boxed{\;1+\iota_{t+1} = \frac{F_{K,t+1}+\left(1-\delta\right)Q_{t+1}}{Q_t}\;}
\label{eq:aq:iota}
\end{align}
Depreciation is valued at $Q_{t+1}$, not at $1$, because replacing worn capital draws material on exactly the same terms as new accumulation does (Remark~\ref{rem:ad:q}). With $\tau^K\equiv0$ this collapses to the familiar $1+\iota_{t+1}=1-\delta+F_{K,t+1}$.

\smalltitle{Euler equation.}
\begin{align}
C_t^{-\sigma} &= \beta\left(1+\iota_{t+1}\right)C_{t+1}^{-\sigma}.
\label{eq:aq:euler}
\end{align}

\smalltitle{Recycling block.} With $Z_t\equiv F_{R,t}+s^R_t$ the private marginal value of a unit of recycled material and $\nu_t\equiv P^W_t$ the private marginal cost of a unit of recyclable waste input, the two firm conditions~\eqref{eq:ad:foc-recycling-b} and~\eqref{eq:ad:foc-waste-use} are
\begin{subequations}\label{eq:aq:recycling-block}
\begin{align}
a'_t\,Z_t &= F_{K,t}
&&\text{(capital in recycling)},
\label{eq:aq:recycling-capital}
\\
a_t\left(1-\varepsilon_t\right)Z_t &= \nu_t
&&\text{(waste input)}.
\label{eq:aq:recycling-input}
\end{align}
\end{subequations}
By~\eqref{eq:aq:recycling-algebra} the ratio of~\eqref{eq:aq:recycling-input} to~\eqref{eq:aq:recycling-capital} is exactly $x_t^2$, so the block collapses to two closed-form statements:
\begin{align}
x_t &= \sqrt{\frac{\nu_t}{F_{K,t}}},
&
\sqrt{F_{K,t}}+\sqrt{\nu_t} &= \sqrt{a^m_tZ_t}.
\label{eq:aq:recycling-closed-form}
\end{align}

\smalltitle{Waste treatment.} From~\eqref{eq:ad:foc-treatment} and~\eqref{eq:aq:CD-derivatives}, the treatment share is closed-form in the price of waste,
\begin{align}
\varpi_t &= \min\left\{1,\ \left[\frac{\left(P^W_t+\tau^W_t\right)\left(1+g^T\right)^t}{\bar c^T}\right]^{1/\varepsilon^T}\right\},
\label{eq:aq:treatment}
\end{align}
where the upper bound is imposed explicitly because $c^T_t$ does not satisfy an Inada condition at $\varpi=1$ (Remark~\ref{rem:ac:varpi}). The zero-profit condition determining the collection subsidy is
\begin{align}
s^W_t &= c^c_t+c^T_t(\varpi_t)+\tau^W_t\left(1-\varpi_t\right)-P^W_t\varpi_t.
\label{eq:aq:zero-profit}
\end{align}

\smalltitle{Extraction and exploration.} Let
\begin{align}
\xi_t &\equiv P^S_t+C^N_{N,t}+\tau^N_t
\label{eq:aq:xi}
\end{align}
denote the full private marginal cost of a unit of virgin material. The extraction condition $F_{R,t}=\xi_t$ together with the production function~\eqref{eq:aq:production} gives a closed form for the materials input net of the physical floor,
\begin{align}
R_t-\bar R &= \left[\frac{\left(1-\alpha\right)A_t\left(K^Y_t\right)^{\alpha}}{\xi_t}\right]^{1/\alpha},
&
Y_t &= A_t\left(K^Y_t\right)^{\alpha}\left(R_t-\bar R\right)^{1-\alpha},
&
N_t &= R_t-a_tW^R_t.
\label{eq:aq:materials-closed-form}
\end{align}
The exploration condition $C^D_{D,t}+P^X_t=P^S_t$ contains no period-$t$ flow, so it inverts into a closed form for the \emph{state} of accumulated discoveries,
\begin{align}
X_t &= \left[\frac{\left(1+\varepsilon^X\right)\left(P^S_t-P^X_t\right)}{c^D_t}\right]^{1/\left(1+\varepsilon^X\right)}
\qquad\text{if}\quad P^S_t>P^X_t,
\label{eq:aq:exploration-closed-form}
\end{align}
with $D_{t-1}=X_t-X_{t-1}$, and $D_{t-1}=0$ otherwise. See Remark~\ref{rem:aq:linear-costs}.

\smalltitle{Shadow prices.} The discrete-time counterparts of~\eqref{eq:ad:PS-PX} are the backward recursions
\begin{align}
P^S_t &= \frac{P^S_{t+1}-C^N_{S,t+1}}{1+\iota_{t+1}},
&
P^X_t &= \frac{P^X_{t+1}+C^D_{X,t+1}}{1+\iota_{t+1}}.
\label{eq:aq:PS-PX}
\end{align}

\smalltitle{Materials and waste.} Market clearing and the ledger:
\begin{align}
W^R_t &= \varpi_tW_t,
\qquad K^R_t = x_tW^R_t,
\qquad
W_t = \frac{N_t-\Delta\Phi_t}{1-a_t\varpi_t}\quad\text{from~\eqref{eq:aq:waste}}.
\label{eq:aq:clearing}
\end{align}

\smalltitle{Summary of the system.} Given policy paths and the initial stocks, equations~\eqref{eq:aq:K-motion}--\eqref{eq:aq:P-motion}, \eqref{eq:aq:euler}, \eqref{eq:aq:recycling-closed-form}, \eqref{eq:aq:treatment}, \eqref{eq:aq:materials-closed-form}, \eqref{eq:aq:exploration-closed-form}, \eqref{eq:aq:PS-PX} and~\eqref{eq:aq:clearing} determine the path of the economy. After the closed-form substitutions the residual unknowns per period are
\begin{align*}
\left\{C_t,\ P^W_t,\ x_t,\ W_t,\ D_t,\ P^S_t,\ P^X_t\right\},
\end{align*}
seven per period, of which $C_t$, $P^S_t$ and $P^X_t$ are forward-looking and the rest are static given the states. $\Omega_t$ and $\phi^Y_t$ are computed afterwards from~\eqref{eq:aq:Omega} wherever they are needed, not solved for jointly; $\phi^K_t=\Phi_t/K_t$ is simply read off the predetermined state $\Phi_t$. The system is a two-point boundary value problem: the stocks are pinned down at $t=0$ and the three forward-looking variables at the terminal date. Variant~\ref{var:aq:finite} states the terminal conditions used in place of the transversality conditions; Part~\ref{part:num} describes how the system is stacked and solved.

\subsection{Regimes and corner conditions}
\label{subsec:aq:regimes}

The model is a complementarity problem, not a smooth system, because of three inequality restrictions:
\begin{align}
K^R_t\ge 0 \quad\perp\quad
&a^m_t\left[F_{R,t}+s^R_t\right]\le F_{K,t},
\label{eq:aq:corner-KR}
\\
W^R_t\ge 0,\ P^W_t\ge 0 \quad\perp\quad
&\text{eq.~\eqref{eq:aq:recycling-input}},
\label{eq:aq:corner-WR}
\\
\varpi_t\in[0,1] \quad\perp\quad
&\text{eq.~\eqref{eq:aq:treatment}}.
\label{eq:aq:corner-varpi}
\end{align}
They generate the three regimes of Section~\ref{subsec:ac:system}:
\begin{enumerate}[label=(\roman*),leftmargin=2.4em]
\item \emph{Linear stage.} $K^R_t=W^R_t=0$, $P^W_t=0$, $a_t=0$, $R_t=N_t$, and $W_t=N_t-\Delta\Phi_t$. If in addition $\tau^W_t=0$ then~\eqref{eq:aq:treatment} gives $\varpi_t=0$ and $s^W_t=c^c_t$: all waste is collected and deposited untreated. The whole recycling block drops out and the model reduces to Variant~\ref{var:aq:linear} below.
\item \emph{Semi-linear stage.} $K^R_t=0$ but $\varpi_t>0$, which under laissez-faire requires $P^W_t>0$ and is therefore impossible; under regulation it is sustained by $\tau^W_t>0$ alone.
\item \emph{Circular stage.} $K^R_t>0$, $P^W_t>0$, and~\eqref{eq:aq:recycling-closed-form} holds with equality.
\end{enumerate}
The date at which~\eqref{eq:aq:corner-KR} first fails --- the \emph{transition date} to the circular economy --- is the central object of the quantitative exercise. Under optimal policy $F_{R,t}+s^R_t=B_t=P^S_t+C^N_{N,t}+d^Wp^P_t$ by~\eqref{eq:ac:B-reduced-B1}, so the transition condition reads
\begin{align}
a^m_t\left[P^S_t+C^N_{N,t}+d^Wp^P_t\right] &> F_{K,t},
\label{eq:aq:transition}
\end{align}
a comparison of the resource rent, extraction cost and pollution damage that a marginal unit of recycling capital saves against the marginal product it forgoes in production. No waste-intensity parameter enters it.

\subsection{Optimal policy}
\label{subsec:aq:optimal-policy}

Implementing the Pigouvian instruments of Proposition~\ref{prop:ad:implementation} requires the planner's shadow prices $p^P_t$, $p^W_t$ and $p^\Phi_t$, computed alongside the market allocation, plus $\Omega_t$ if the rates are to be reported in interpretable units. Throughout, $B_t\equiv P^S_t+C^N_{N,t}+d^Wp^P_t$ and the social return replaces the private one in~\eqref{eq:aq:iota} via $Q_t=q_t$.

\smalltitle{Pollution.}
\begin{align}
p^P_t &= \frac{MSC^P_{t+1}+\left(1-\widehat\theta_{t+1}\right)p^P_{t+1}}{1+\iota_{t+1}},
&
MSC^P_t &= \frac{\psi_tP_t^{\phi}}{\zeta_t}+\gamma\frac{Y_t}{P_t},
&
\zeta_t &= C_t^{-\sigma}.
\label{eq:aq:pP}
\end{align}
$p^P_t$ is the present value at $t$ of a unit of $P_{t+1}$, matching the one-period lag in~\eqref{eq:aq:P-motion}.

\smalltitle{Waste.} The ledger price is explicit, from~\eqref{eq:ac:pW-explicit-B1}:
\begin{align}
p^W_t &= c^c_t+c^T_t(\varpi_t)+p^P_t\left(1-\varpi_t+d^W\varpi_t\right)-\varpi_t\,a_t\left(1-\varepsilon_t\right)B_t.
\label{eq:aq:pW}
\end{align}
No simultaneous solve is required: given $p^P_t$, $P^S_t$ and the period's allocation, \eqref{eq:aq:pW} returns $p^W_t$ directly. Its sign is not restricted (Remark~\ref{rem:aq:pW-sign}).

\smalltitle{Embodied capital material.} $p^\Phi_t$ is the present value at $t$ of the waste liability attached to one unit of $\Phi_{t+1}$: that unit releases $\delta$ as waste at $t+1$ and leaves $\left(1-\delta\right)$ units of $\Phi_{t+2}$ behind. Hence
\begin{align}
p^\Phi_t &= \frac{\delta\,p^W_{t+1}+\left(1-\delta\right)p^\Phi_{t+1}}{1+\iota_{t+1}},
&
q_t &= 1-\phi^I_t\left(p^W_t-p^\Phi_t\right).
\label{eq:aq:pPhi-discrete}
\end{align}
Because $\delta+\left(1-\delta\right)<1+\iota_{t+1}$ whenever $\iota_{t+1}>0$, a constant $p^W$ gives $p^\Phi_t<p^W_t$ and hence $q_t<1$; the general statement is~\eqref{eq:ac:q-explicit}.

\smalltitle{Instrument levels.} The discrete-time counterparts of~\eqref{eq:ad:implement} are
\begin{align}
\tau^N_t &= p^W_t,
&
s^R_t &= d^Wp^P_t-p^W_t,
&
\tau^W_t &= p^P_t\left(1-d^W\right),
&
\tau^K_t &= -\phi^I_t\left(p^W_t-p^\Phi_t\right).
\label{eq:aq:instruments}
\end{align}
Under these, $Q_t=q_t$ exactly, so the market's Euler equation~\eqref{eq:aq:euler} reduces to the planner's own saving condition; this is what makes the capital-formation margin, and not only the waste- and pollution-related margins, decentralize correctly.

\smalltitle{Practical consequence.} Under optimal policy the market and planner allocations coincide, so in practice one solves the \emph{planner} system --- which is square and has no market price $P^W_t$, the shadow price $\nu_t=dc^T_t/d\varpi-p^P_t\left(1-d^W\right)$ taking its place in~\eqref{eq:aq:recycling-closed-form} --- and then reads off the instruments from~\eqref{eq:aq:instruments}. Solving the market system with the instruments imposed is then a validation exercise: the two allocations must agree.

Three further checks are informative, and are more useful than a residual test on the materials balance, which now holds by construction:
\begin{enumerate}[label=(\arabic*),leftmargin=2.4em]
\item $F_{R,t}-C^N_{N,t}-p^W_t=P^S_t$ and $B_t=F_{R,t}-p^W_t+d^Wp^P_t$ must agree with $B_t=P^S_t+C^N_{N,t}+d^Wp^P_t$ to machine precision. This is the cancellation of Section~\ref{subsec:ac:reductions} and is the sharpest available test of the accounting.
\item $P^W_t=a_t\left(1-\varepsilon_t\right)B_t$ must be strictly positive throughout the circular stage.
\item $\Omega_t$ from~\eqref{eq:aq:Omega} must stay positive and must not drift to implausible values. A negative $\Omega_t$ means $\phi^I_tG_t>N_t+R^R_t$ --- more mass embodied in new capital than entered the economy --- which is the genuine feasibility restriction identified in Remark~\ref{rem:ac:ledger} and signals a calibration inconsistency in $\bar\phi^I$ or $\varsigma$, not a coding error.
\end{enumerate}

\subsection{Variants}
\label{subsec:aq:variants}

The following restricted versions of the model are useful for building up and testing the implementation. Each is obtained by a parameter restriction on the full model, so the same code path can serve all of them.

\begin{enumerate}[label=\textbf{V\arabic*.},leftmargin=3.2em,itemsep=0.4em]

\item\label{var:aq:linear} \emph{Linear-economy restriction.} Impose $K^R_t=W^R_t=\varpi_t=0$ for all $t$. Then $a_t=0$, $R_t=N_t$, $K^Y_t=K_t$, $C^W_t=c^c_tW_t$, $W_t=N_t-\Delta\Phi_t$ with no circularity multiplier, and $P_{t+1}=\left[1-\theta(P_t)\right]P_t+W_t$. This restriction leaves $\Phi_t$'s own law of motion~\eqref{eq:aq:Phi-motion-discrete} untouched; combine with Variant~\ref{var:aq:constant-phi} to drop it as well. The entire recycling block, the price $P^W_t$ and the treatment margin disappear; the remaining unknowns are $\left\{C_t,N_t,D_t,W_t,P^S_t,P^X_t\right\}$, and $p^W_t=c^c_t+p^P_t$ in closed form. This variant is the natural starting point: it is a standard Ramsey model with exhaustible resources plus a pollution stock, and it produces the initial conditions from which the full model departs.

\item\label{var:aq:no-feedback} \emph{No productivity feedback from pollution.} Set $\gamma=0$, so $A_t=\bar A\left(1+g^A\right)^t$. Under exogenous policy the pollution stock then influences nothing: \eqref{eq:aq:P-motion} becomes a pure forward recursion that can be run \emph{after} the allocation is solved. Under optimal policy the decomposition breaks, because $p^P_t$ still feeds back through $\tau^W_t$, $s^R_t$ and $p^W_t$, but the second term of $MSC^P_t$ in~\eqref{eq:aq:pP} drops.

\item\label{var:aq:no-resource} \emph{Exogenous resource base.} Fix $S_t=\bar S$ and $X_t=\bar X$, or equivalently set $\varepsilon^S=\varepsilon^X=0$ and drop~\eqref{eq:aq:SX-motion}. Extraction cost becomes a constant marginal cost $c^N_t/\left(1+\varepsilon^S\right)$, exploration disappears, and $P^S_t=P^X_t=0$. Note this makes $B_t=C^N_{N,t}+d^Wp^P_t$, so the transition date~\eqref{eq:aq:transition} is driven by pollution and extraction cost alone.

\item\label{var:aq:no-floor} \emph{No materials floor.} Set $\bar R=0$. Then $F_{R,t}=(1-\alpha)Y_t/R_t$ and the extraction condition becomes the constant-share relation $\xi_tR_t=(1-\alpha)Y_t$.

\item\label{var:aq:no-progress} \emph{No technical progress.} Set $g^A=g^R=g^N=g^D=g^c=g^T=g^\psi=g^i_\omega=0$ for $i\in\{Y,N,D,C,W\}$, and, under closing rule~(a) of Section~\ref{subsec:aq:phiI}, $g^\phi=0$ as well (under rule~(b) there is no separate growth rate to zero out, since $\phi^I_t=\varsigma\Omega_t$ inherits whatever path $\Omega_t$ takes). Note that this list has no $g^\Omega$: there is no such parameter in this document, since $\Omega_t$ is derived rather than calibrated. Its path under this restriction is still generally non-constant, because it depends on the allocation; the same is true of $\phi^K_t$, which continues to evolve via~\eqref{eq:aq:Phi-motion-discrete} even with $g^\phi=0$, converging to $\bar\phi^I$ only asymptotically.

\item\label{var:aq:finite} \emph{Finite horizon.} Truncate at $T$ and replace the transversality conditions~\eqref{eq:ac:transversality} by terminal conditions. Treating period-$T$ values as constant thereafter gives the perpetuity terminal values
\begin{align}
P^S_T &= \frac{-C^N_{S,T}}{\iota_T},
&
P^X_T &= \frac{C^D_{X,T}}{\iota_T},
&
p^P_T &= \frac{MSC^P_T}{\iota_T+\widehat\theta_T},
&
p^\Phi_T &= \frac{\delta\,p^W_T}{\iota_T+\delta},
&
K_{T+1} &= K_T,
\label{eq:aq:terminal}
\end{align}
where the last condition imposes zero net investment in the terminal period; $\Phi_{T+1}=\Phi_T$ closes $\Phi$'s own path alongside it, which by~\eqref{eq:aq:Phi-motion-discrete} requires $\phi^I_T\delta K_T=\delta\Phi_T$, i.e.\ $\phi^K_T=\phi^I_T$ --- the vintage composition has converged. All quantitative results are computed on the finite-horizon version, with $T$ chosen large enough that the objects of interest are insensitive to it. That insensitivity should be reported.

\item\label{var:aq:constant-phi} \emph{Constant material intensity.} Impose $\phi^I_t\equiv\phi^K_t\equiv\bar\phi^K$ for all $t$ (equivalently, set $g^\phi=0$ in closing rule~(a) and $\Phi_0=\bar\phi^KK_0$, so~\eqref{eq:aq:Phi-motion-discrete} holds with $\Phi_t\equiv\bar\phi^KK_t$ identically). $\Phi_t$ then drops out as an independent state, $\Delta\Phi_t=\bar\phi^KI_t$, and~\eqref{eq:aq:waste} becomes $W_t=\left(N_t-\bar\phi^KI_t\right)/\left(1-a_t\varpi_t\right)$: only \emph{net} investment stores material, since replacement returns exactly what it draws. The liability $p^\Phi_t$ and the instrument $\tau^K_t=-\bar\phi^K\left(p^W_t-p^\Phi_t\right)$ survive unchanged --- capital still defers waste. Useful as a baseline against which to measure how much the vintage-composition channel actually moves the results.

\end{enumerate}

\subsection{Scenarios}
\label{subsec:aq:scenarios}

The scenarios combine a \emph{policy} configuration (Section~\ref{subsec:ad:implementation}) with a \emph{technical progress} configuration.

\smalltitle{Benchmark: optimal policy and neutral technical progress.}
\begin{align}
g^A=g^R=g^N=g^D=g^c=g^T=g^\psi>0,
\qquad
g^Y_\omega=g^N_\omega=g^D_\omega=g^C_\omega=g^W_\omega=0,
\qquad
g^\phi=0\ \text{(rule~(a))}.
\label{eq:aq:benchmark}
\end{align}
All sectors and the marginal disutility weight $\psi_t$ grow at a common rate, so no sector is favoured by assumption and the transition to the circular economy is driven by capital accumulation, resource depletion and pollution accumulation alone. Under closing rule~(a), $g^\phi=0$ means new-vintage capital carries a constant intensity $\bar\phi^I$, the same neutrality the $g^i_\omega=0$ restrictions impose elsewhere; under rule~(b) there is no separate choice to make. There is no growth rate for $\Omega$ to hold at zero: under this benchmark $\Omega_t$'s path is an \emph{output} of the model, worth reporting alongside the transition dates as a check that it stays within a plausible range. $\phi^K_t$'s path is worth reporting for the same reason.

\smalltitle{Experiments.} Departures from~\eqref{eq:aq:benchmark} of interest:
\begin{enumerate}[label=(\alph*),leftmargin=2.4em]
\item \emph{Sector-specific technical progress under optimal policy.} Vary each of $g^R$, $g^N$, $g^D$, $g^c$, $g^T$ and (under closing rule~(a)) $g^\phi$ one at a time around the benchmark and record the effect on the date of first waste treatment and the date of first recycling. Note that the relative intensities' growth rates $g^i_\omega$ are \emph{not} in this list under rule~(a): by~\eqref{eq:aq:waste} they do not touch the allocation, and varying them moves only the reported $\Omega_t$. Under rule~(b) they do matter, through $\phi^I_t=\varsigma\Omega_t$, which makes the comparison between the two closing rules an experiment in its own right. $g^\phi$ is the natural way to isolate the vintage-composition channel: a strongly positive $g^\phi$ (materials-saving capital-goods technology) shrinks the sink term $\Delta\Phi_t$ in~\eqref{eq:aq:waste} and should, other things equal, raise waste and bring the transition forward.
\item \emph{Laissez-faire with neutral progress.} All four instruments zero, with $s^W_t=c^c_t+c^T_t(\varpi_t)-P^W_t\varpi_t$ from~\eqref{eq:aq:zero-profit}. This answers whether an unregulated market economy reaches the circular stage at all, and if so how late.
\item \emph{Partial regulation with neutral progress.} $\tau^W_t=s^R_t=0$, with $\tau^N_t$ and $\tau^K_t$ at their Pigouvian levels. This isolates the contribution of the waste-treatment and recycling instruments specifically. The complementary experiment --- $\tau^K_t=0$ with the other three at their Pigouvian levels --- isolates the capital-formation channel, and is the natural way to quantify how much the materials-sink property of capital is worth.
\end{enumerate}
Scenarios (b) and (c) require the welfare measure of Section~\ref{subsec:aq:welfare} to be comparable across policies, and (a) requires the optimal-policy block of Section~\ref{subsec:aq:optimal-policy} to be re-solved for each parameterization.

\subsection{Calibration}
\label{subsec:aq:calibration}

\smalltitle{Initial condition.} The economy is assumed to start in an early linear stage in which the environment absorbs all waste generated, so the pollution stock is initially stationary. With $t=0$ the first period,
\begin{align}
P_1 &= \left[1-\theta(P_0)\right]P_0+W_0 = P_0
&&\Longleftrightarrow&&
\bar\theta P_0^{1-\varepsilon^P} = W_0.
\label{eq:aq:calibration-pollution}
\end{align}
Equation~\eqref{eq:aq:calibration-pollution} is one restriction linking $\bar\theta$, $\varepsilon^P$, $P_0$ and the initial waste flow $W_0$; given data on two of them it determines a third. $W_0=N_0-\Delta\Phi_0$ from~\eqref{eq:aq:waste} in the linear stage, not a separately calibrated quantity; it requires $\Phi_0$ as an additional initial condition alongside $K_0,S_0,X_0,P_0$.

\smalltitle{Initial regime.} The assumption that period $0$ lies in the linear stage is itself a restriction on the calibration and must be checked rather than imposed:
\begin{align}
a^m_0\left[F_{R,0}+s^R_0\right] &\le F_{K,0},
&&\text{with }a^m_0=1-b.
\label{eq:aq:calibration-linear}
\end{align}
Since $a^m_0=1-b$, the parameter $b$ controls how far the initial economy is from the recycling threshold and therefore, jointly with $g^R$, the transition date.

\smalltitle{Parameters set directly.} $\alpha$ from the capital income share; $\delta$ and $\beta$ (equivalently $\rho$) from conventional values given the length of a period; $\sigma$ from the intertemporal elasticity of substitution; $\varepsilon^S$, $\varepsilon^X$, $\varepsilon^T$, $\varepsilon^P$, $\gamma$ and $d^W$ from the relevant empirical literatures; the growth rates $g^\cdot$ from the scenario definitions. There is no single material intensity $\phi^K$ to set directly: under closing rule~(a), $\bar\phi^I$ and $g^\phi$ are set from material-flow accounts for capital-goods production specifically, if available; under rule~(b), only $\varsigma$ needs setting.

The \emph{relative} waste intensities $\omega^Y,\omega^N,\omega^D,\omega^C,\omega^W$ occupy an unusual position and should be described as such. Under closing rule~(a) they do not enter the allocation at all (Section~\ref{subsec:aq:materials}); they are recovered-quantity parameters that determine the reported $\Omega_t$ and $\phi^Y_t$, and hence what the model says about material intensities in the data, but they cannot be identified from --- and do not affect --- the transition dates. They should therefore be set from material-flow accounts and then \emph{checked} against the implied $\Omega_t$, rather than calibrated to match a model outcome. Under rule~(b) they do affect the allocation, through $\phi^I_t=\varsigma\Omega_t$, and the distinction between the two rules is worth reporting.

\smalltitle{Parameters solved for.} The scale constants $c^N$, $c^D$, $\bar c^c$ and $\bar c^T$ are set so that the model reproduces realistic shares of the mining and waste industries in output in the base year. The preference parameters $\bar\psi$ and $\phi$ govern the marginal willingness to pay for environmental quality and are the natural targets for a stated- or revealed-preference estimate of the marginal damage of pollution in the base year. The initial stocks $K_0$, $\Phi_0$, $S_0$, $X_0$ and $P_0$ are data, subject to~\eqref{eq:aq:calibration-pollution}; $\Phi_0$ is most naturally reported as $\phi^K_0=\Phi_0/K_0$, an initial average material intensity of the capital stock.

\smalltitle{Nested structure.} Any target that involves an equilibrium object --- the industry shares, the base-year marginal damage --- cannot be evaluated without solving the model, so the final stage of the calibration is a nested fixed-point problem: an outer loop over $\left(c^N,c^D,\bar c^c,\bar c^T,\bar\psi,b\right)$ around an inner solve of the equilibrium system. Part~\ref{part:num} describes the loop.

\subsection{Welfare measurement}
\label{subsec:aq:welfare}

Split lifetime utility~\eqref{eq:aq:utility} into its consumption and pollution components,
\begin{align}
U_0 = U^C_0-U^P_0,
\qquad
U^C_0 = \sum_{t=0}^{\infty}\beta^t\frac{C_t^{1-\sigma}}{1-\sigma},
\qquad
U^P_0 = \sum_{t=0}^{\infty}\beta^t\frac{\psi_t}{1+\phi}P_t^{1+\phi}.
\label{eq:aq:welfare-split}
\end{align}

\smalltitle{The household's lifetime budget constraint.} Let $\Delta_t\equiv\prod_{i=1}^{t}\left(1+\iota_i\right)^{-1}$ with $\Delta_0=1$ be the goods discount factor implied by~\eqref{eq:aq:iota}. Discounting the household budget~\eqref{eq:ad:budget} in the form $C_t=\Pi^H_t+T_t-Q_tK_{t+1}+Q_t\left(1-\delta\right)K_t$ and telescoping the capital terms using $\Delta_tQ_t=\Delta_{t+1}\left[F_{K,t+1}+\left(1-\delta\right)Q_{t+1}\right]$ gives
\begin{align}
\sum_{t=0}^{\infty}\Delta_tC_t &= H_0+A_0,
\label{eq:aq:lifetime-budget}
\end{align}
where
\begin{align}
A_0 &\equiv \left[F_{K,0}+\left(1-\delta\right)Q_0\right]K_0,
\label{eq:aq:A0}
\\
H_0 &\equiv \sum_{t=0}^{\infty}\Delta_t\left[\Pi^H_t+T_t-F_{K,t}K_t\right]
= \sum_{t=0}^{\infty}\Delta_t\left[Y_t-F_{K,t}K_t-C^N_t-C^D_t-\left(c^c_t+c^T_t\right)W_t+\tau^K_tG_t\right].
\label{eq:aq:H0}
\end{align}
$A_0$ is the value at $t=0$ of the predetermined capital stock --- its current marginal product plus its resale value --- and $H_0$ is the present value of pure profits, including resource rents and the transfer of the capital-formation subsidy. Writing initial wealth this way resolves the boundary problem that arises if one instead discounts $K_0$ back one period: no lagged instrument $\tau^K_{-1}$ appears anywhere.

\smalltitle{Consumption growth and the closed form.} Iterating~\eqref{eq:aq:euler} forward,
\begin{align}
w_t &\equiv \frac{C_t}{C_0}=\left[\beta^t\prod_{i=1}^{t}\left(1+\iota_i\right)\right]^{1/\sigma},
\label{eq:aq:weights}
\end{align}
so that with the intertemporal price index $P^u_0$ and the utility index $Z_0$,
\begin{align}
P^u_0 &\equiv \sum_{t=0}^{\infty}\Delta_tw_t,
&
Z_0 &\equiv \sum_{t=0}^{\infty}\beta^tw_t^{1-\sigma},
&
C_0 &= \frac{H_0+A_0}{P^u_0},
\label{eq:aq:indices}
\end{align}
lifetime consumption utility admits the closed form
\begin{align}
U^C_0 &= \frac{Z_0}{1-\sigma}\left(\frac{H_0+A_0}{P^u_0}\right)^{1-\sigma}.
\label{eq:aq:UC}
\end{align}
Every occurrence of the discount factor $\prod\left(1+r_i\right)^{-1}$ in a specification without a capital-formation instrument is replaced here, term by term, by $\prod\left(1+\iota_i\right)^{-1}$ with $\iota$ from~\eqref{eq:aq:iota}, since that is the rate at which the household actually trades resources across periods once capital is priced at $Q$ rather than at $1$. With $\tau^K\equiv0$ everything reduces to the familiar expressions.

\smalltitle{Compensating variation.} Let starred objects denote values along the first-best path and unstarred objects values along the actual path. The welfare cost of deviating from the optimum is $\Delta W_0$ defined by
\begin{align}
\frac{Z_0}{1-\sigma}\left(\frac{H_0+A_0+\Delta W_0}{P^u_0}\right)^{1-\sigma}
&= U^{C*}_0+U^P_0-U^{P*}_0,
\notag \\
\Longleftrightarrow\qquad
\Delta W_0 &= P^u_0\left[\left(\frac{1-\sigma}{Z_0}\right)\left(U^{C*}_0+U^P_0-U^{P*}_0\right)\right]^{1/\left(1-\sigma\right)}-\left(H_0+A_0\right).
\label{eq:aq:compensating-variation}
\end{align}
$Z_0$, $P^u_0$, $H_0$, $A_0$ and $U^P_0$ are evaluated on the \emph{actual} path, so~\eqref{eq:aq:compensating-variation} requires both paths to be solved and stored.

\subsection{Technical remarks}
\label{subsec:aq:remarks}

\begin{remark}[What is given up by \textbf{(B1)}]\label{rem:aq:scope}
Setting $\omega^{N,S}=\omega^{D,S}=0$ is what makes the ledger collapse to the closed form~\eqref{eq:aq:waste}, and with it what makes $\Omega_t$ allocation-irrelevant under closing rule~(a). Section~\ref{subsec:aq:B1} lists the five things it buys; this remark records what it costs, now that Sections~\ref{sec:ac}--\ref{sec:ad} state the general case explicitly and the comparison can be made precisely rather than by assertion.

The first cost is physical: overburden and gangue displaced by extraction, and any analogous byproduct of exploration, are no longer distinguished from waste generated by the final-good expenditure on those activities. Empirically this is not a small omission --- displaced material typically dominates ore tonnage by an order of magnitude --- and the model has nothing to say about it under \textbf{(B1)}.

The second is structural and is the more important. With either intensity positive, $\Omega_t$ re-enters the first-order conditions, $\left(W_t,\Omega_t\right)$ must be solved jointly with the rest of the period's block rather than in two clean steps, and the extraction margin acquires both a displacement term and a scale term. The relative-intensity calibration, currently a reporting choice, would become a determinant of the transition date. Three points are worth stating precisely, because an earlier version of this remark understated the change:
\begin{enumerate}[label=(\arabic*),leftmargin=2.4em,itemsep=0.2em]
\item The extra term in the extraction condition is not merely $-\Omega_t\omega^{N,S}p^W_t$. Because $\Omega_t$ is itself endogenous, the general condition~\eqref{eq:ac:foc-extraction} also carries the scale term $-p^W_t\kappa_t$ and the dilution terms in $\pi_t$. The same is true of every other margin.
\item The ledger is a \emph{quadratic} in $W_t$, \eqref{eq:ac:ledger-quadratic}, not an implicit equation requiring iteration --- so the cost is one closed form replaced by another, not a loss of closed form. The unique positive root is what the solver would take.
\item The instrument set grows from four to seven, and $p^W_t$ no longer cancels from $B_t$. The latter is the more consequential: it means the transition date~\eqref{eq:aq:transition} stops being independent of the waste-intensity parameters, which is one of the cleanest results the \textbf{(B1)} model delivers.
\end{enumerate}
Reintroducing the split is therefore cheaper computationally than it first appears and more expensive interpretively. It changes the character of the calibration exercise, not only its cost --- and, before it is attempted, the specification question raised in Remark~\ref{rem:ac:dilution} should be settled, since the absolute-intensity alternative removes most of the added complexity while keeping all of the added physical content.
\end{remark}

\begin{remark}[The sign of $p^W_t$ and what the solver must tolerate]\label{rem:aq:pW-sign}
By~\eqref{eq:aq:pW}, $p^W_t<0$ whenever $\varpi_ta_t(1-\varepsilon_t)B_t$ exceeds $c^c_t+c^T_t+p^P_t\left(1-\varpi_t+d^W\varpi_t\right)$, and $B_t$ contains $P^S_t$, which grows without bound as reserves deplete. Three consequences for the implementation. First, $\tau^N_t=p^W_t$ turns into an extraction \emph{subsidy}, correcting the fact that extraction stocks the recycling loop --- a benefit external to the extractor. Second, $q_t>1$ and $\tau^K_t>0$: when waste is a resource, storing material in capital is socially costly rather than valuable, and the two flip together, which is a useful internal check. Third, no bound on $p^W_t$, $\tau^N_t$ or $\tau^K_t$ should be imposed in the solver, and the complementarity conditions of Section~\ref{subsec:aq:regimes} should be verified to behave sensibly through the sign change. Whether the regime is actually reached in the calibrated model is an empirical question worth reporting either way.
\end{remark}

\begin{remark}[Linear cost functions and the exploration margin]\label{rem:aq:linear-costs}
Because~\eqref{eq:aq:extraction-cost} and~\eqref{eq:aq:discovery-cost} are linear in $N_t$ and $D_t$, marginal extraction and discovery costs are independent of the current flow. For extraction this is harmless: the condition $F_{R,t}=\xi_t$ still pins down $N_t$, because $F_{R,t}$ is decreasing in $R_t$. For exploration it is not: $C^D_{D,t}$ depends only on the state $X_t$, so the exploration condition determines $X_t$, as in~\eqref{eq:aq:exploration-closed-form}, and hence $D_{t-1}$; the exploration margin is intertemporal rather than static. $D_t$ should be recovered from the path of $X$, not solved for directly, and the model has a bang-bang flavour on this margin --- $D_{t-1}$ jumps whenever $P^S_t-P^X_t$ moves. If that turns out to be quantitatively unattractive, the fix is to make $C^D$ strictly convex in $D$, e.g.\ $C^D_t=\frac{c^D_t}{1+\varepsilon^X}\frac{D_t^{1+\varepsilon^D}}{1+\varepsilon^D}X_t^{1+\varepsilon^X}$ with $\varepsilon^D>0$, at the cost of losing~\eqref{eq:aq:exploration-closed-form}.
\end{remark}

\begin{remark}[Scaling and numerical conditioning]\label{rem:aq:scaling}
With $g^A>0$ the model has no stationary steady state, and $Y_t$, $K_t$ and $C_t$ grow without bound while $S_t$ falls towards zero. Five practical points. First, the equilibrium system should be solved in detrended or log variables; the natural deflator is $\left(1+g^A\right)^{t/(1-\alpha)}$ for the goods block, but the resource and pollution blocks do not share it, so the blocks should be scaled separately. Second, $P^S_t$ and $P^X_t$ grow at rates unrelated to output and are the most likely source of ill-conditioning. Third, $\varpi_t$, $x_t$, $a_t$ and $q_t$ are bounded objects and should be solved for directly rather than recovered from ratios of unbounded ones. Fourth, $\Phi_t$ inherits $K_t$'s unbounded growth and should be scaled the same way, but the derived ratio $\phi^K_t=\Phi_t/K_t$ is bounded and, under either closing rule for $\phi^I_t$, should converge to a well-behaved path rather than drift --- persistent divergence between $\phi^K_t$ and $\phi^I_t$ is economically meaningful (a capital stock that never catches up to the current vintage) but should be checked against the calibrated $g^\phi$ or $\varsigma$ before being read as a result. Fifth, $\Omega_t$ from~\eqref{eq:aq:Omega} is a ratio of two growing objects; under closing rule~(a) an $\Omega_t$ that drifts towards zero, explodes, or turns negative is a diagnostic about the calibration of the $\omega^i$'s and $\bar\phi^I$, and cannot be a coding error in the allocation block, since the allocation block never sees it.
\end{remark}

\begin{remark}[The circularity multiplier and the period length]\label{rem:aq:period}
The factor $\left(1-a_t\varpi_t\right)^{-1}$ in~\eqref{eq:aq:waste} has material extracted at the start of period $t$ being recycled and re-used within period $t$ (Remark~\ref{rem:ac:recirculation}). This is exact in continuous time but becomes a substantive assumption once the period length is fixed, and it is not innocuous quantitatively: as $a_t\varpi_t\to1$ the multiplier diverges, so $R_t$ and $W_t$ both become large relative to $N_t$ late in the transition. The alternative --- lagging recycled input, $R^R_t=a_{t-1}\varpi_{t-1}W_{t-1}$ --- removes the multiplier at the cost of an extra state and destroys the closed form~\eqref{eq:aq:materials-closed-form}, since $R_t$ would then be predetermined. Reporting $R_t/N_t$ and $W_t/N_t$ alongside the transition dates makes the size of the effect visible, and the sensitivity of the results to the period length is worth checking directly.
\end{remark}


\clearpage
\part{Numerical}\label{part:num}

This part contains a description of the numerical approaches used to solve the quantitative model of Part~\ref{part:models}. None of the numerical strategy depends on the process-level accounting as such, only on the closed-form structure of the equilibrium system worked out in Section~\ref{subsec:aq:system}. That structure rests on the baseline restriction \textbf{(B1)} of Section~\ref{subsec:aq:B1} and on nothing else; Remark~\ref{rem:aq:scope} records what would change without it, which is less than one might expect --- the ledger remains closed-form, as the unique positive root of a quadratic.



% \printbibliography[
% heading=bibintoc,
% title={References},
% category=cited]

%%%%%%%%%%%%%%%%%%%%%%%%%%%%%%%%%%%%%%%%%%%%%%%%%%%%%%
%%                      Appendices                  %%
%%%%%%%%%%%%%%%%%%%%%%%%%%%%%%%%%%%%%%%%%%%%%%%%%%%%%%

% \begin{appendices}
% \end{appendices}


\end{document}
